%% bare_jrnl.tex
%% V1.4b
%% 2015/08/26
%% by Michael Shell
%% see http://www.michaelshell.org/
%% for current contact information.
%%
%% This is a skeleton file demonstrating the use of IEEEtran.cls
%% (requires IEEEtran.cls version 1.8b or later) with an IEEE
%% journal paper.
%%
%% Support sites:
%% http://www.michaelshell.org/tex/ieeetran/
%% http://www.ctan.org/pkg/ieeetran
%% and
%% http://www.ieee.org/

%%*************************************************************************
%% Legal Notice:
%% This code is offered as-is without any warranty either expressed or
%% implied; without even the implied warranty of MERCHANTABILITY or
%% FITNESS FOR A PARTICULAR PURPOSE! 
%% User assumes all risk.
%% In no event shall the IEEE or any contributor to this code be liable for
%% any damages or losses, including, but not limited to, incidental,
%% consequential, or any other damages, resulting from the use or misuse
%% of any information contained here.
%%
%% All comments are the opinions of their respective authors and are not
%% necessarily endorsed by the IEEE.
%%
%% This work is distributed under the LaTeX Project Public License (LPPL)
%% ( http://www.latex-project.org/ ) version 1.3, and may be freely used,
%% distributed and modified. A copy of the LPPL, version 1.3, is included
%% in the base LaTeX documentation of all distributions of LaTeX released
%% 2003/12/01 or later.
%% Retain all contribution notices and credits.
%% ** Modified files should be clearly indicated as such, including  **
%% ** renaming them and changing author support contact information. **
%%*************************************************************************


% *** Authors should verify (and, if needed, correct) their LaTeX system  ***
% *** with the testflow diagnostic prior to trusting their LaTeX platform ***
% *** with production work. The IEEE's font choices and paper sizes can   ***
% *** trigger bugs that do not appear when using other class files.       ***                          ***
% The testflow support page is at:
% http://www.michaelshell.org/tex/testflow/



\documentclass[journal]{IEEEtran}
%
% If IEEEtran.cls has not been installed into the LaTeX system files,
% manually specify the path to it like:
% \documentclass[journal]{../sty/IEEEtran}





% Some very useful LaTeX packages include:
% (uncomment the ones you want to load)


% *** MISC UTILITY PACKAGES ***
%
%\usepackage{ifpdf}
% Heiko Oberdiek's ifpdf.sty is very useful if you need conditional
% compilation based on whether the output is pdf or dvi.
% usage:
% \ifpdf
%   % pdf code
% \else
%   % dvi code
% \fi
% The latest version of ifpdf.sty can be obtained from:
% http://www.ctan.org/pkg/ifpdf
% Also, note that IEEEtran.cls V1.7 and later provides a builtin
% \ifCLASSINFOpdf conditional that works the same way.
% When switching from latex to pdflatex and vice-versa, the compiler may
% have to be run twice to clear warning/error messages.






% *** CITATION PACKAGES ***
%
%\usepackage{cite}
% cite.sty was written by Donald Arseneau
% V1.6 and later of IEEEtran pre-defines the format of the cite.sty package
% \cite{} output to follow that of the IEEE. Loading the cite package will
% result in citation numbers being automatically sorted and properly
% "compressed/ranged". e.g., [1], [9], [2], [7], [5], [6] without using
% cite.sty will become [1], [2], [5]--[7], [9] using cite.sty. cite.sty's
% \cite will automatically add leading space, if needed. Use cite.sty's
% noadjust option (cite.sty V3.8 and later) if you want to turn this off
% such as if a citation ever needs to be enclosed in parenthesis.
% cite.sty is already installed on most LaTeX systems. Be sure and use
% version 5.0 (2009-03-20) and later if using hyperref.sty.
% The latest version can be obtained at:
% http://www.ctan.org/pkg/cite
% The documentation is contained in the cite.sty file itself.






% *** GRAPHICS RELATED PACKAGES ***
%
\ifCLASSINFOpdf
  % \usepackage[pdftex]{graphicx}
  % declare the path(s) where your graphic files are
  % \graphicspath{{../pdf/}{../jpeg/}}
  % and their extensions so you won't have to specify these with
  % every instance of \includegraphics
  % \DeclareGraphicsExtensions{.pdf,.jpeg,.png}
\else
  % or other class option (dvipsone, dvipdf, if not using dvips). graphicx
  % will default to the driver specified in the system graphics.cfg if no
  % driver is specified.
  % \usepackage[dvips]{graphicx}
  % declare the path(s) where your graphic files are
  % \graphicspath{{../eps/}}
  % and their extensions so you won't have to specify these with
  % every instance of \includegraphics
  % \DeclareGraphicsExtensions{.eps}
\fi
% graphicx was written by David Carlisle and Sebastian Rahtz. It is
% required if you want graphics, photos, etc. graphicx.sty is already
% installed on most LaTeX systems. The latest version and documentation
% can be obtained at: 
% http://www.ctan.org/pkg/graphicx
% Another good source of documentation is "Using Imported Graphics in
% LaTeX2e" by Keith Reckdahl which can be found at:
% http://www.ctan.org/pkg/epslatex
%
% latex, and pdflatex in dvi mode, support graphics in encapsulated
% postscript (.eps) format. pdflatex in pdf mode supports graphics
% in .pdf, .jpeg, .png and .mps (metapost) formats. Users should ensure
% that all non-photo figures use a vector format (.eps, .pdf, .mps) and
% not a bitmapped formats (.jpeg, .png). The IEEE frowns on bitmapped formats
% which can result in "jaggedy"/blurry rendering of lines and letters as
% well as large increases in file sizes.
%
% You can find documentation about the pdfTeX application at:
% http://www.tug.org/applications/pdftex





% *** MATH PACKAGES ***
%
%\usepackage{amsmath}
% A popular package from the American Mathematical Society that provides
% many useful and powerful commands for dealing with mathematics.
%
% Note that the amsmath package sets \interdisplaylinepenalty to 10000
% thus preventing page breaks from occurring within multiline equations. Use:
%\interdisplaylinepenalty=2500
% after loading amsmath to restore such page breaks as IEEEtran.cls normally
% does. amsmath.sty is already installed on most LaTeX systems. The latest
% version and documentation can be obtained at:
% http://www.ctan.org/pkg/amsmath





% *** SPECIALIZED LIST PACKAGES ***
%
%\usepackage{algorithmic}
% algorithmic.sty was written by Peter Williams and Rogerio Brito.
% This package provides an algorithmic environment fo describing algorithms.
% You can use the algorithmic environment in-text or within a figure
% environment to provide for a floating algorithm. Do NOT use the algorithm
% floating environment provided by algorithm.sty (by the same authors) or
% algorithm2e.sty (by Christophe Fiorio) as the IEEE does not use dedicated
% algorithm float types and packages that provide these will not provide
% correct IEEE style captions. The latest version and documentation of
% algorithmic.sty can be obtained at:
% http://www.ctan.org/pkg/algorithms
% Also of interest may be the (relatively newer and more customizable)
% algorithmicx.sty package by Szasz Janos:
% http://www.ctan.org/pkg/algorithmicx




% *** ALIGNMENT PACKAGES ***
%
%\usepackage{array}
% Frank Mittelbach's and David Carlisle's array.sty patches and improves
% the standard LaTeX2e array and tabular environments to provide better
% appearance and additional user controls. As the default LaTeX2e table
% generation code is lacking to the point of almost being broken with
% respect to the quality of the end results, all users are strongly
% advised to use an enhanced (at the very least that provided by array.sty)
% set of table tools. array.sty is already installed on most systems. The
% latest version and documentation can be obtained at:
% http://www.ctan.org/pkg/array


% IEEEtran contains the IEEEeqnarray family of commands that can be used to
% generate multiline equations as well as matrices, tables, etc., of high
% quality.




% *** SUBFIGURE PACKAGES ***
%\ifCLASSOPTIONcompsoc
%  \usepackage[caption=false,font=normalsize,labelfont=sf,textfont=sf]{subfig}
%\else
%  \usepackage[caption=false,font=footnotesize]{subfig}
%\fi
% subfig.sty, written by Steven Douglas Cochran, is the modern replacement
% for subfigure.sty, the latter of which is no longer maintained and is
% incompatible with some LaTeX packages including fixltx2e. However,
% subfig.sty requires and automatically loads Axel Sommerfeldt's caption.sty
% which will override IEEEtran.cls' handling of captions and this will result
% in non-IEEE style figure/table captions. To prevent this problem, be sure
% and invoke subfig.sty's "caption=false" package option (available since
% subfig.sty version 1.3, 2005/06/28) as this is will preserve IEEEtran.cls
% handling of captions.
% Note that the Computer Society format requires a larger sans serif font
% than the serif footnote size font used in traditional IEEE formatting
% and thus the need to invoke different subfig.sty package options depending
% on whether compsoc mode has been enabled.
%
% The latest version and documentation of subfig.sty can be obtained at:
% http://www.ctan.org/pkg/subfig




% *** FLOAT PACKAGES ***
%
%\usepackage{fixltx2e}
% fixltx2e, the successor to the earlier fix2col.sty, was written by
% Frank Mittelbach and David Carlisle. This package corrects a few problems
% in the LaTeX2e kernel, the most notable of which is that in current
% LaTeX2e releases, the ordering of single and double column floats is not
% guaranteed to be preserved. Thus, an unpatched LaTeX2e can allow a
% single column figure to be placed prior to an earlier double column
% figure.
% Be aware that LaTeX2e kernels dated 2015 and later have fixltx2e.sty's
% corrections already built into the system in which case a warning will
% be issued if an attempt is made to load fixltx2e.sty as it is no longer
% needed.
% The latest version and documentation can be found at:
% http://www.ctan.org/pkg/fixltx2e


%\usepackage{stfloats}
% stfloats.sty was written by Sigitas Tolusis. This package gives LaTeX2e
% the ability to do double column floats at the bottom of the page as well
% as the top. (e.g., "\begin{figure*}[!b]" is not normally possible in
% LaTeX2e). It also provides a command:
%\fnbelowfloat
% to enable the placement of footnotes below bottom floats (the standard
% LaTeX2e kernel puts them above bottom floats). This is an invasive package
% which rewrites many portions of the LaTeX2e float routines. It may not work
% with other packages that modify the LaTeX2e float routines. The latest
% version and documentation can be obtained at:
% http://www.ctan.org/pkg/stfloats
% Do not use the stfloats baselinefloat ability as the IEEE does not allow
% \baselineskip to stretch. Authors submitting work to the IEEE should note
% that the IEEE rarely uses double column equations and that authors should try
% to avoid such use. Do not be tempted to use the cuted.sty or midfloat.sty
% packages (also by Sigitas Tolusis) as the IEEE does not format its papers in
% such ways.
% Do not attempt to use stfloats with fixltx2e as they are incompatible.
% Instead, use Morten Hogholm'a dblfloatfix which combines the features
% of both fixltx2e and stfloats:
%
% \usepackage{dblfloatfix}
% The latest version can be found at:
% http://www.ctan.org/pkg/dblfloatfix




%\ifCLASSOPTIONcaptionsoff
%  \usepackage[nomarkers]{endfloat}
% \let\MYoriglatexcaption\caption
% \renewcommand{\caption}[2][\relax]{\MYoriglatexcaption[#2]{#2}}
%\fi
% endfloat.sty was written by James Darrell McCauley, Jeff Goldberg and 
% Axel Sommerfeldt. This package may be useful when used in conjunction with 
% IEEEtran.cls'  captionsoff option. Some IEEE journals/societies require that
% submissions have lists of figures/tables at the end of the paper and that
% figures/tables without any captions are placed on a page by themselves at
% the end of the document. If needed, the draftcls IEEEtran class option or
% \CLASSINPUTbaselinestretch interface can be used to increase the line
% spacing as well. Be sure and use the nomarkers option of endfloat to
% prevent endfloat from "marking" where the figures would have been placed
% in the text. The two hack lines of code above are a slight modification of
% that suggested by in the endfloat docs (section 8.4.1) to ensure that
% the full captions always appear in the list of figures/tables - even if
% the user used the short optional argument of \caption[]{}.
% IEEE papers do not typically make use of \caption[]'s optional argument,
% so this should not be an issue. A similar trick can be used to disable
% captions of packages such as subfig.sty that lack options to turn off
% the subcaptions:
% For subfig.sty:
% \let\MYorigsubfloat\subfloat
% \renewcommand{\subfloat}[2][\relax]{\MYorigsubfloat[]{#2}}
% However, the above trick will not work if both optional arguments of
% the \subfloat command are used. Furthermore, there needs to be a
% description of each subfigure *somewhere* and endfloat does not add
% subfigure captions to its list of figures. Thus, the best approach is to
% avoid the use of subfigure captions (many IEEE journals avoid them anyway)
% and instead reference/explain all the subfigures within the main caption.
% The latest version of endfloat.sty and its documentation can obtained at:
% http://www.ctan.org/pkg/endfloat
%
% The IEEEtran \ifCLASSOPTIONcaptionsoff conditional can also be used
% later in the document, say, to conditionally put the References on a 
% page by themselves.
\usepackage{graphicx}
\usepackage{amsfonts}       % blackboard math symbols
\usepackage{amsmath}
\usepackage{amsthm}
%\usepackage{amssymb}
\usepackage{soul}
\usepackage{color}


% \usepackage{natbib}  % DO NOT CHANGE THIS AND DO NOT ADD ANY OPTIONS TO IT
\usepackage{hyperref}       % hyperlinks
\usepackage{url}            % simple URL typesetting
\usepackage{booktabs}       % professional-quality tables
\urlstyle{same}

\usepackage{nicefrac}       % compact symbols for 1/2, etc.
\usepackage{microtype}      % microtypography
\usepackage{comment}
\usepackage{bm}
\usepackage{multirow}
\usepackage{enumerate}
\usepackage{lineno}
\usepackage{makecell}

\usepackage{algorithm}
\usepackage[noend]{algpseudocode}

\newtheorem{definition}{Definition}
\newtheorem{theorem}{Theorem}

% basic
\newcommand{\RR}{\mathbb{R}}                                     % real numbers
\newcommand{\CC}{\mathbb{C}}                                     % complex numbers
\newcommand{\diag}{{\rm diag}}
\newcommand{\Diag}{{\rm Diag}\;}
\newcommand{\trace}{{\rm trace}}
%\newcommand{\sin}{{\rm sin}}
%\newcommand{\cos}{{\rm cos}}
\newcommand{\sech}{{\rm sech}}
\newcommand{\Prox}{{\rm Prox}}
\newcommand{\dist}{{\rm dist}}
\newcommand{\vectorize}{{\rm vec}}
%\newcommand{\tanh}{{\rm tanh}}
\newcommand{\bi}{\begin{itemize}}
\newcommand{\ei}{\end{itemize}}
\newcommand{\ba}{\begin{array}}
\newcommand{\ea}{\end{array}}
\newcommand{\DefinedAs}[0]{\mathrel{\mathop:}=}
\newcommand{\non}{\nonumber}
\DeclareMathOperator*{\argmin}{arg\,min}
\DeclareMathOperator*{\argmax}{arg\,max}

% model
\usepackage{xspace}
\newcommand{\model}{{TGV-CRN}\xspace}
\newcommand{\modelnospace}{{TGV-CRN}}
%\newcommand{\model}{\texttt{GVCRN}\xspace}
%\newcommand{\modelnospace}{\texttt{GVCRN}}

% theorem
\newtheorem{assumption}{Assumption}[section]
\newcommand{\indep}{\perp\!\!\!\!\perp} 

% imitation 
\newcommand{\pivec}{\vec{\pi}}
\newcommand{\ovec}{\vec{o}}
\newcommand{\svec}{\vec{s}}
\newcommand{\avec}{\vec{a}}
\newcommand{\E}{\mathbb{E}}

% VRNN
\newcommand{\qt}[1]{\left<#1\right>}
\newcommand{\ql}[1]{\left[#1\right]}
\newcommand{\hess}{\mathbf{H}}
\newcommand{\jacob}{\mathbf{J}}
% \newcommand{\hl}{HL}
\newcommand{\cost}{\mathcal{L}}
\newcommand{\lout}{\mathbf{r}}
\newcommand{\louti}{r}
\newcommand{\outi}{y}
\newcommand{\out}{\mathbf{y}}
\newcommand{\gauss}{\mathbf{G_N}}
\newcommand{\eye}{\mathbf{I}}
\newcommand{\softmax}{\phi}
\newcommand{\targ}{\mathbf{t}}
\newcommand{\metric}{\mathbf{G}}
\newcommand{\sample}{\mathbf{z}}

\newcommand{\eq}[1]{\begin{align}#1\end{align}}
\newcommand{\eqn}[1]{\begin{align*}#1\end{align*}}
\newcommand{\eps}{\epsilon}
\newcommand{\kl}{\texttt{KL}}
\newcommand{\Expect}{\operatorname{\mathbb{E}}}

\newcommand{\bmx}[0]{\begin{bmatrix}}
\newcommand{\emx}[0]{\end{bmatrix}}
\newcommand{\qexp}[1]{\left<#1\right>}
\newcommand{\vect}[1]{\mathbf{#1}}
\newcommand{\vects}[1]{\boldsymbol{#1}}
\newcommand{\matr}[1]{\mathbf{#1}}
\newcommand{\var}[0]{\operatorname{Var}}
\newcommand{\std}[0]{\operatorname{std}}
\newcommand{\cov}[0]{\operatorname{Cov}}
\newcommand{\matrs}[1]{\boldsymbol{#1}}
\newcommand{\va}[0]{\vect{a}}
\newcommand{\vb}[0]{\vect{b}}
\newcommand{\vc}[0]{\vect{c}}
\newcommand{\vh}[0]{\vect{h}}
% \newcommand{\vv}[0]{\vect{v}}
\newcommand{\vx}[0]{\vect{x}}
\newcommand{\vz}[0]{\vect{z}}
\newcommand{\vw}[0]{\vect{w}}
\newcommand{\vs}[0]{\vect{s}}
\newcommand{\vf}[0]{\vect{f}}
\newcommand{\vi}[0]{\vect{i}}
\newcommand{\vo}[0]{\vect{o}}
\newcommand{\vy}[0]{\vect{y}}
\newcommand{\vg}[0]{\vect{g}}
\newcommand{\vm}[0]{\vect{m}}
\newcommand{\vu}[0]{\vect{u}}
\newcommand{\vL}[0]{\vect{L}}
\newcommand{\vr}[0]{\vect{r}}
\newcommand{\mW}[0]{\matr{W}}
\newcommand{\mG}[0]{\matr{G}}
\newcommand{\mX}[0]{\matr{X}}
\newcommand{\mQ}[0]{\matr{Q}}
\newcommand{\mU}[0]{\matr{U}}
\newcommand{\mV}[0]{\matr{V}}
\newcommand{\mA}{\matr{A}}
\newcommand{\mD}{\matr{D}}
\newcommand{\mS}{\matr{S}}
\newcommand{\mI}{\matr{I}}
\newcommand{\td}[0]{\text{d}}
\newcommand{\TT}[0]{\vects{\theta}}
\newcommand{\vepsilon}[0]{\vects{\epsilon}}
\newcommand{\vsigma}[0]{\vects{\sigma}}
\newcommand{\valpha}[0]{\vects{\alpha}}
\newcommand{\vmu}[0]{\vects{\mu}}
\newcommand{\vzero}[0]{\vect{0}}
\newcommand{\tf}[0]{\text{m}}
\newcommand{\tdf}[0]{\text{dm}}
\newcommand{\grad}[0]{\nabla}
\newcommand{\alert}[1]{\textcolor{black}{#1}}
\newcommand{\N}[0]{\mathcal{N}}
\newcommand{\LL}[0]{\mathcal{L}}
\newcommand{\MM}[0]{\mathcal{M}}
\newcommand{\NN}[0]{\mathcal{N}}
\newcommand{\HH}[0]{\mathcal{H}}
\newcommand{\II}[0]{\mathbb{I}}
\newcommand{\Scal}[0]{\mathcal{S}}
\newcommand{\sigmoid}{\sigma}
\newcommand{\enabla}[0]{\ensuremath{%
    \overset{\raisebox{-0.3ex}[0.5ex][0ex]{%
    \ensuremath{\scriptscriptstyle e}}}{\nabla}}}
\newcommand{\enhnabla}[0]{\nabla_{\hspace{-0.5mm}e}\,}
\newcommand{\tred}[1]{\textcolor{black}{#1}}
\newcommand{\todo}[1]{{\Large\textcolor{black}{#1}}}
\newcommand{\done}[1]{{\Large\textcolor{green}{#1}}}
\newcommand{\dd}[1]{\ensuremath{\mbox{d}#1}}


% *** PDF, URL AND HYPERLINK PACKAGES ***
%
%\usepackage{url}
% url.sty was written by Donald Arseneau. It provides better support for
% handling and breaking URLs. url.sty is already installed on most LaTeX
% systems. The latest version and documentation can be obtained at:
% http://www.ctan.org/pkg/url
% Basically, \url{my_url_here}.




% *** Do not adjust lengths that control margins, column widths, etc. ***
% *** Do not use packages that alter fonts (such as pslatex).         ***
% There should be no need to do such things with IEEEtran.cls V1.6 and later.
% (Unless specifically asked to do so by the journal or conference you plan
% to submit to, of course. )


% correct bad hyphenation here
\hyphenation{op-tical net-works semi-conduc-tor}


\begin{document}
%
% paper title
% Titles are generally capitalized except for words such as a, an, and, as,
% at, but, by, for, in, nor, of, on, or, the, to and up, which are usually
% not capitalized unless they are the first or last word of the title.
% Linebreaks \\ can be used within to get better formatting as desired.
% Do not put math or special symbols in the title.
\title{Estimating counterfactual treatment outcomes over time in complex multiagent scenarios}
%
%
% author names and IEEE memberships
% note positions of commas and nonbreaking spaces ( ~ ) LaTeX will not break
% a structure at a ~ so this keeps an author's name from being broken across
% two lines.
% use \thanks{} to gain access to the first footnote area
% a separate \thanks must be used for each paragraph as LaTeX2e's \thanks
% was not built to handle multiple paragraphs
%

\author{Keisuke Fujii,~\IEEEmembership{Member,~IEEE,}
        Koh Takeuchi, 
        Atsushi Kuribayashi,
        Naoya Takeishi,\\
        Yoshinobu Kawahara,
        and~Kazuya Takeda,~\IEEEmembership{Senior~Member,~IEEE}% <-this % stops a space
\thanks{K. Fujii is with Graduate School of Informatics at Nagoya University, JAPAN and RIKEN Center for Advanced Intelligence Project, JAPAN. e-mail: fujii@i.nagoya-u.ac.jp.}% <-this % stops a space
\thanks{K. Takeuchi is with Graduate School of Informatics at Kyoto University and RIKEN.}
\thanks{A. Kuribayashi and K. Takeda are with Graduate School of Informatics at Nagoya University.}% <-this % stops a space
\thanks{N. Takeishi is with Graduate School of Engineering at the University of Tokyo and RIKEN Center for Advanced Intelligence Project.}
\thanks{Y. Kawahara is with Graduate School of Information Science and Technology at Osaka University and RIKEN Center for Advanced Intelligence Project.}}
% \thanks{Manuscript received April 19, 2005; revised August 26, 2015.}}

% note the % following the last \IEEEmembership and also \thanks - 
% these prevent an unwanted space from occurring between the last author name
% and the end of the author line. i.e., if you had this:
% 
% \author{....lastname \thanks{...} \thanks{...} }
%                     ^------------^------------^----Do not want these spaces!
%
% a space would be appended to the last name and could cause every name on that
% line to be shifted left slightly. This is one of those "LaTeX things". For
% instance, "\textbf{A} \textbf{B}" will typeset as "A B" not "AB". To get
% "AB" then you have to do: "\textbf{A}\textbf{B}"
% \thanks is no different in this regard, so shield the last } of each \thanks
% that ends a line with a % and do not let a space in before the next \thanks.
% Spaces after \IEEEmembership other than the last one are OK (and needed) as
% you are supposed to have spaces between the names. For what it is worth,
% this is a minor point as most people would not even notice if the said evil
% space somehow managed to creep in.



% The paper headers %  Transactions on Neural Networks and Learning Systems /  of Machine Learning
\markboth{Submitted to IEEE TNNLS, Special Issue: Information Theoretic Methods for the Generalization, Robustness and Interpretability...}% Journal of \LaTeX\ Class Files,~Vol.~14, No.~8, August~2015}%
{Fujii \MakeLowercase{\textit{et al.}}: Estimating counterfactual treatment outcomes over time in multiagent}
% The only time the second header will appear is for the odd numbered pages
% after the title page when using the twoside option.
% 
% *** Note that you probably will NOT want to include the author's ***
% *** name in the headers of peer review papers.                   ***
% You can use \ifCLASSOPTIONpeerreview for conditional compilation here if
% you desire.




% If you want to put a publisher's ID mark on the page you can do it like
% this:
%\IEEEpubid{0000--0000/00\$00.00~\copyright~2015 IEEE}
% Remember, if you use this you must call \IEEEpubidadjcol in the second
% column for its text to clear the IEEEpubid mark.



% use for special paper notices
%\IEEEspecialpapernotice{(Invited Paper)}




% make the title area
\maketitle

% As a general rule, do not put math, special symbols or citations
% in the abstract or keywords.
\begin{abstract}
Evaluation of intervention in a multiagent system, e.g., when humans should intervene in autonomous driving systems and when a player should pass to teammates for a good shot, is challenging in various engineering and scientific fields. 
Estimating the individual treatment effect (ITE) using counterfactual long-term prediction is practical to evaluate such interventions.
However, most of the conventional frameworks did not consider the time-varying complex structure of multiagent relationships and covariate counterfactual prediction. 
This may lead to erroneous assessments of ITE and difficulty in interpretation.
Here we propose an interpretable, counterfactual recurrent network in multiagent systems to estimate the effect of the intervention.
Our model leverages graph variational recurrent neural networks and theory-based computation with domain knowledge for the ITE estimation framework based on long-term prediction of multiagent covariates and outcomes, which can confirm the circumstances under which the intervention is effective.   
On simulated models of an automated vehicle and biological agents with time-varying confounders, we show that our methods achieved lower estimation errors in counterfactual covariates and the most effective treatment timing than the baselines. Furthermore, using real basketball data, our methods performed realistic counterfactual predictions and evaluated the counterfactual passes in shot scenarios. 
\end{abstract}

% Note that keywords are not normally used for peerreview papers.
\begin{IEEEkeywords}
multiagent modeling, causal inference, deep generative model, trajectory data, autonomous vehicle, sports.
\end{IEEEkeywords}






% For peer review papers, you can put extra information on the cover
% page as needed:
% \ifCLASSOPTIONpeerreview
% \begin{center} \bfseries EDICS Category: 3-BBND \end{center}
% \fi
%
% For peerreview papers, this IEEEtran command inserts a page break and
% creates the second title. It will be ignored for other modes.
\IEEEpeerreviewmaketitle



\section{Introduction}
% The very first letter is a 2 line initial drop letter followed
% by the rest of the first word in caps.
% 
% form to use if the first word consists of a single letter:
% \IEEEPARstart{A}{demo} file is ....
% 
% form to use if you need the single drop letter followed by
% normal text (unknown if ever used by the IEEE):
% \IEEEPARstart{A}{}demo file is ....
% 
% Some journals put the first two words in caps:
% \IEEEPARstart{T}{his demo} file is ....
% 
% Here we have the typical use of a "T" for an initial drop letter
% and "HIS" in caps to complete the first word.
Evaluation of intervention in a real-world multiagent system is a fundamental problem in a variety of engineering and scientific fields. 
For example, a human driver in an autonomous vehicle, a player in team sports, and an experimenter on animals intervene in multiagent systems to obtain desirable results (e.g., safe driving, a good shot, and specific behavior, respectively) as shown in Fig. \ref{fig:intervention}.  
In these processes and complex interactions between agents, it is often difficult to estimate the intervention (or treatment) effect to compare the outcomes with and without interventions. 
There have been many methods to estimate the individual treatment effect (ITE), which evaluates the causal effect of treatment strategies on some important outcomes at the individual level in various fields (e.g., \cite{glass2013causal,baum2015causal,wang2015robust}).
In particular, some work has been proposed for dealing with time-varying \cite{lim2018forecasting, bica2020estimating} and hidden confounders \cite{bica2020time,liu2020estimating,ma2021deconfounding}. 


However, most of the conventional frameworks did not consider the time-varying complex structures of multiagent relationships and counterfactual covariate predictions, which may lead to erroneous assessments of ITE and difficulty of interpretation. 
The structures of multiagent relationships include bottom-up ones based on interactions between local agents (often represented as a graph), and top-down ones described by global statistics and/or theories with domain knowledge \cite{Vicsek95,Couzin02}.
Since real-world multiagent systems do not usually have explicit governing equations, both top-down and bottom-up approaches can be simultaneously required for the modeling \cite{Fujii18,Fujii20}.
That is, in our problem, we need to model multiagent systems in both data-driven and theory-based approaches. 
In addition, for extracting insights from the ITE estimation results, the circumstances where intervention is effective must be analyzed.   
Therefore, plausible and interpretable modeling of time-varying multiagent systems is required for estimating ITE.  

In this paper, we propose a novel causal inference framework called Theory-based and Graph Variational Counterfactual Recurrent Network (\modelnospace), which estimates the effect of intervention in multiagent systems in interpretable ways.
The causal graph of the problem setting is illustrated in Fig. \ref{fig:causalgraph}, where the hidden confounders at a particular time stamp not only have causal relations to the observed variables at the same time stamp but also are causally affected by the hidden confounders from previous time stamps.
% both the hidden confounders and the treatment assignment from previous time stamps \cite{bica2020time}.
To model the hidden confounders and other causal relationships from data for the ITE estimation framework, we leverage graph variational recurrent neural networks (GVRNNs) \cite{Yeh19} to represent local agent interactions and theory-based computation for global properties of the multiagent behaviors based on domain knowledge.
This framework is based on the long-term prediction of multiagent covariates and outcomes can confirm under what circumstances the intervention is effective.   

Our general motivation is to estimate ITE over time in complex multiagent scenarios.
Specifically, in team sports, decision-making skills (e.g., whether a player with the ball should perform a pass or shot) are important. However, we cannot observe both patterns as data (i.e., performing a pass \textit{and} a shot) in the same situation. 
For autonomous driving or animal behavioral science, in real-world scenarios, we cannot obtain both data (i.e., with and without intervention) in the same situation (note that in the numerical experiment, we used synthetic data including both cases). For animal behavioral science, similarly, scientific experiments can hardly be performed with all possible intervention timing during movements and sometimes include the bias of the experimenters. 

In summary, our main contributions are as follows.
(1) We proposed a novel counterfactual recurrent network called \modelnospace, which estimates the effect of intervention in \textcolor{black}{\textit{multiagent systems in interpretable ways}, compared with the previous counterfactual recurrent networks dealing with time-varying treatment \cite{lim2018forecasting, bica2020estimating,bica2020time,liu2020estimating,ma2021deconfounding}.} 
(2) Methodologically, for the ITE estimation framework based on long-term prediction of multiagent covariates and outcomes, our model leverages GVRNN to represent local agent interactions and theory-based computation, \textcolor{black}{which was not considered in the above previous work. This framework} can confirm under what circumstances the intervention is effective.
(3) In experiments using two simulated models of an automated vehicle and biological agents, we show that our methods achieved lower errors in estimating counterfactual covariates and the most effective treatment timing than the baselines. Furthermore, using real basketball data, our methods performed realistic counterfactual predictions.
\textcolor{black}{All of these subjects moved interactively in multiagent systems, which are not dealt with in the previous work.}
We extend our previous short paper \cite{fujii2022estimating} by adding theoretical background, experimental results of a synthetic Boid dataset and a real-world basketball dataset, and a sensitivity analysis using the CARLA dataset.   

The remainder of this paper is organized as follows. First, in Section \ref{sec:problem}, we describe our problem definition. Next, we describe our methods in Section \ref{sec:proposed}.
We overview the related works in Section \ref{sec:related}, present experimental results in Section \ref{sec:experiments}, and conclude this paper in Section \ref{sec:conclusion}.

\begin{figure}[t]
\centering
\includegraphics[width=1\linewidth]{fig/Fig_intervention.png}
\caption[]{The illustrations of our problems. Interventions in (A) an autonomous vehicle simulation, (B) a biological agent simulation, and (C) a real basketball are shown.
In (A) and (C), a single agent (A: an ego-vehicle and C: a ball player) is intervened whereas multiple agents (all boids) are intervened in (B).
\textcolor{black}{The motivations and variable definitions are described in the Introduction and Background sections.
In short, we aim to perform long-term counterfactual prediction of outcomes and covariates from the past covariates and intervention (or treatment assignment).}
In (A), our approach can test autonomous driving software including human interventions without creating the same situations for controlled trials (as real-world scenarios).
\textcolor{black}{The outcome is a safe driving distance.}
In (B), our approach has the potential to estimate the effect of an experimenter’s interventions on multi-animal behaviors, which improves the efficiency of experimental procedures for observing desired movements.
\textcolor{black}{The outcome is the angular velocity of multiple agents.}
In (C), our approach can estimate the effect of the selection of passes in basketball shot scenarios, thus we can evaluate the decision-making skills in this situation (e.g., during a game). \textcolor{black}{The outcome is the effectiveness of an attack.}

}

\label{fig:intervention}
\end{figure}

\section{Background} % Problem definition}
\label{sec:problem}
In this section, we first give definitions of the notations used throughout the paper, present the assumptions of our methods for estimating ITE, \textcolor{black}{and then introduce VRNNs (variational recurrent neural networks) \cite{Chung15} and GNN (graph neural network) \cite{Kipf18} we used}.

\subsection{Preliminary}
The multiagent observational data is denoted as the following $X_t, A_t, Y_t$ at time stamp $t$. % , C \in\mathcal{X}_t
Examples of these variables are shown in Fig. \ref{fig:intervention}\textcolor{black}{.}
Let $X_t$ be the time-dependent covariates of the observational data such that $X_t=\{x^{(1)}_{t},...,x^{(n)}_{t}\}$, where the $x^{(i)}_{t}$ denotes the covariates for $i$-th multiagent sample with $K$ agents, and $n$ denotes the number of samples. %, and $\mathcal{X}_t$ denotes the time-dependent feature space. 
The relationships between agents are represented by the theory-based computation and GVRNN described in Section \ref{sec:proposed}. 
Although some related papers \cite{bica2020estimating,liu2020estimating} consider the static covariates $C$, % \in\mathcal{C}
which do not change over time, here we do not explicitly consider $C$ because we can easily formulate and implement our methods to add $C$ by conditioning. 
The treatment (or intervention) assignments are denoted as $A_{t}=\{a^{(1)}_{t},a^{(2)}_{t},...,a^{(n)}_{t}\}$, where $a^{(i)}_{t}$ denotes the treatments assigned in the $i$-th sample. % \in\mathcal{A}
We consider $a^{i}_{t}\in\{0, 1\}$, where $1$ is considered treated whereas $0$ is the control (i.e., a binary treatment setting), and we estimate the effect of the treatment assigned at time stamp $t$ on the outcomes $Y_{t+1}=\{y^{(1)}_{t+1},y^{(2)}_{t+1},...,y^{(n)}_{t+1}\}$ at time stamp $t+1$.
% , where $\tau$ is the prediction window. \in\mathcal{Y}
Note that in observational data, a multiagent sample can only belong to one group (i.e., either a treated or control group), thus the outcome from the other group is always missing and referred to as counterfactual.
To represent the historical sequential data before time stamp $t$, we use the notation $\overline{X}_{t}=\{X_{1},X_{2},...,X_{t-1}\}$ to denote the history of covariates observed before time stamp $t$, and $\overline{A}_{t}$ refers to the history of treatment assignments. Combining all covariates and treatments, we define $\mathcal{H}^{(i)}_{t}=\{\overline{x}^{(i)}_{t},\overline{a}^{(i)}_{t}\}$ as all the historical data collected before time stamp $t$. % \cup\{c^{(i)}\}
% The observational data for $i$-th sample can be represented using the notations defined above as: $\mathcal{D}^{(i)}=\{x^{(i)}_{t},a^{(i)}_{t}\}^{T}_{t=1}\cup\{c^{(i)},y^{(i)}_{a,T+\tau}\}$. 
% The detailed description of the notations can be referred to Table \ref{tab:notations}.

% In this paper, we build our framework upon the well-adopted potential outcome framework \cite{Neyman23, Rubin05} [22, 28] and its variation that considers the time-varying treatment assignments when estimating the causal effect on the outcomes. 
We adopt the potential outcomes framework (e.g., \cite{rubin1978bayesian}) and extended by \cite{robins2009estimation} to account for time-varying treatments.
The potential outcome $y^{(i)}_{a_t=a,t+1}$ of the $i$-th sample given the historical treatment can be formulated as $y^{(i)}_{a_t=a,t+1}=\mathbb{E}[y|x^{(i)}_{t},\mathcal{H}^{(i)}_t,a_t=a]$, where $a=\{0,1\}$. 
Then the ITE on the temporal observational data is defined as:
\begin{small}
\begin{equation}
    \tau^{(i)}_t=\mathbb{E}[y^{(i)}_{a_t =1,t+1}-y^{(i)}_{a_t =0,t+1}|x^{(i)}_{t},\mathcal{H}^{(i)}_{t}].
\end{equation}
\end{small}
Here, the observed outcome $y^{(i)}_{a_t=a,t+1}$ under treatment $a$ is called factual outcome, while the unobserved one $y^{(i)}_{a_t=1-a,t+1}$ is the counterfactual outcome.
% In observational data, only the factual outcomes are available, while the counterfactual outcomes can never be observed. 


\begin{figure}[t]
\centering
\includegraphics[width=0.8\linewidth]{fig/Fig_causalgraph.png}
\caption[]{The illustration of causal graphs for our problem.
We denote $X_t, Z_t,  A_t, Y_{t+1}$ as the dynamic covariates, representations of hidden confounders, treatment assignment, and outcomes, respectively. % C, static covariates, 
The black lines indicate the causal relations.
The hidden confounders $Z_{t+1}$ usually affect the treatment assignment $A_{t+1}$, the outcome $Y_{t+2}$, and the covariate $X_{t}$. To infer $Z_{t+1}$, we can leverage the observational data $X_{t+1}$ and previous hidden counfounders $Z_{t}$.
%The hidden confounders $Z_{t+1}$ usually affect the treatment assignment $A_{t+1}$ the outcome $Y_{t+2}$. To infer $Z_{t+1}$, we can leverage the observational data $X_{t+1}$, previous hidden confounders $Z_{t}$ and treatment assignment $A_{t}$. % and $C$
}

\label{fig:causalgraph}
\end{figure}

\subsection{Assumptions}
Our estimation of ITE is based on the following standard assumptions \cite{robins2000marginal, lim2018forecasting,hernan2010causal}, and we further extend the assumptions in our scenario to include time-varying observational data.

\begin{assumption}[Consistency]\label{as:consistency}
The potential outcome under treatment history $\overline{A}$ is equal to the observed outcome if the actual treatment history is $\overline{A}$.
\end{assumption}

\begin{assumption}[Positivity]\label{as:positivity}
For any sample $i$, if the probability  $p(\overline{a}^{(i)}_{t-1},\overline{x}^{(i)}_{t})\neq0$, then the probability of receiving treatment $0$ or $1$ is positive, i.e., % c^{(i)}
$0<p(\overline{a}^{(i)}_{t},\overline{x}^{(i)}_{t} )<1$, for all $\overline{a}^{(i)}_{t}$. % ,c^{(i)}  
\end{assumption}
% \mathbb{P}
Assumption \ref{as:positivity} means that, for each time $t$, each treatment has a non-zero probability of being assigned. 
% Most existing works (e.g.,  \cite{hill2011bayesian,shalit2017estimating}) rely on the strong ignorability assumption \cite{rosenbaum1983central} as described below, assuming that observed features are enough to eliminate the confounding bias, i.e., no hidden confounders exist.
Besides these two assumptions, much of the existing work is based on the \textit{strong ignorability} assumption as follows:

\begin{definition}[Sequential Strong Ignorability]\label{as:ignorability}
Given the observed historical covariates $\overline{x}^{(i)}_t$ % and static covariates $c^{(i)}$ 
of the $i$-th sample, the potential outcome variables $\{y^{(i)}_{a_t=0,t+1}, y^{(i)}_{a_t=1,t+1}\}$ are independent of the treatment assignment, i.e., $\{y^{(i)}_{a_t=0,t+1}, y^{(i)}_{a_t=1,t+1}\}\indep  a^{(i)}_{t}|\overline{x}^{(i)}_t$. % ,c^{(i)}
\end{definition}

Definition \ref{as:ignorability} means that there are no hidden confounders, i.e., all covariates affecting both the treatment assignment and the outcomes are present in the observational dataset. 
However, this condition is difficult to be guaranteed in practice especially in real-world observational data (in other words, it is not testable in practice \cite{robins2000marginal,  pearl2009causality}).
In this paper, we relax such a strict assumption by acknowledging potential hidden confounders. Our proposed methods can learn the representations of the hidden confounders and eliminate the bias between the treatment assignments and outcomes at each time stamp.

In our approach, the learned representations (denoted by $Z_{t}=\{z^{(1)}_{t},z^{(2)}_{t},...,z^{(n)}_{t}\}$) can be leveraged to infer the unobserved confounders and act as substitutes of hidden confounders. % \in\mathcal{Z}
That is, we extend the strong ignorability assumption by considering the existence of hidden confounders $Z_{t}$, which influence the treatment assignment $A_t$ and potential outcomes $Y_{t+1}$. 
Given the hidden confounders $Z_{t}$, the potential outcome variables are independent of the treatment assignment at each time stamp. 
We aim to learn the representations of hidden confounders $Z_{t}$ for bias elimination based on the following assumptions \cite{ma2021deconfounding}: 

\begin{assumption}[Existence of Hidden Confounders]\label{as:hiddenconfounders}
(i) The hidden confounders may not be accessible, but the covariates are correlated with the hidden confounders, and can be considered as proxy variables, and
(ii) hidden confounders at each time stamp are also influenced by the hidden confounders and treatment assignments from previous time stamps.
\end{assumption}
% Given the hidden confounders $Z_{t}$, the potential outcome variables are independent of the treatment assignment at each time stamp. 
%We aim to learn the representations of hidden confounders for bias elimination based on the following assumptions \cite{ma2021deconfounding}: 
%(i) The hidden confounders may not be accessible, but the covariates are correlated with the hidden confounders, and can be considered as proxy variables, and
%(ii) hidden confounders at each time stamp are also influenced by the hidden confounders and treatment assignments from previous time stamps.

Based on the premise, we study the identification of ITE: 
\setcounter{theorem}{0}
\begin{theorem}[Identification of ITE]\label{thm:identification}
If we recover $p(z_t^{(i)}|x_t^{(i)},\mathcal{H}_t^{(i)})$ and $p(y_{t+1}^{(i)}|z_t^{(i)},a_t^{(i)})$, then the proposed methods can recover the ITE under the causal graph in Fig. \ref{fig:causalgraph}.
\end{theorem}
\noindent We provide a proof in Appendix \ref{app:proof}.
For simplicity, the sample superscript ($i$) will be omitted unless explicitly needed.


\begin{figure*}[t]
\centering
\includegraphics[width=1\textwidth]{fig/Fig_model.png}
\caption[]{The illustration of \model. 
\textcolor{black}{(A) \model aims to estimate ITE based on long-term prediction of multiagent covariates and outcomes while visualizing the long-term future covariate prediction. \model leverages GVRNN to represent local agent interactions and theory-based functions for covariate and outcome prediction, which can confirm under what circumstances the intervention is effective. Specifically, (B) the training and inference processes of GNN encoder, prior, and decoder are illustrated.}
At each time stamp, the model takes the current covariates and treatment assignments as input to learn representations of the hidden confounders via GRNNs and GNN encoders. 
Then, via theory-based computations, the GNN decoders, and MLPs (multi-layer perceptron), the model predicts time-varying covariates, a potential outcome, and a treatment. \textcolor{black}{We also use the gradient reversal layer before the treatment classifier to ensure the confounder representation distribution of the treated and that of the controlled are similar at the group level.}
}

\label{fig:model}
\end{figure*}

\textcolor{black}{
\subsection{Variational recurrent and graph neural networks}
}
\textcolor{black}{Here we explain VRNN \cite{Chung15} and GNN \cite{Kipf18} used in the following section \ref{sec:proposed}. }

\textcolor{black}{
\noindent{\bf{VRNN.}}
Let $x_{\leq T} = \{ x_1, \dots, x_T \}$ denote a sequence of variables of length $T$. 
% Here we consider single-agent modeling for simplicity.
The goal of sequential generative modeling is to learn the distribution over sequential data $\mathcal{D}$ consisting of multiple demonstrations. 
% We assume that all sequences have the same length $T$, but in general, this does not need to be the case.
A common approach to model the trajectory is to factorize the joint distribution and then maximize the log-likelihood 
%
%\begin{align}
$\theta^* = \argmax_{\theta} \sum_{x_{\leq T} \in \mathcal D} \sum_{t=1}^T \log p_{\theta} (x_t | x_{<t}),$ % = \argmax_{\theta} \sum_{x_{\leq T} \in \mathcal D} \log p_{\theta} (x_{\leq T}) 
%label{eq:condprobs}
%end{align}
%
where $\theta$ denotes the learnable parameters of models such as RNNs.
% recurrent neural networks (RNNs).
However, RNNs with simple output distributions often struggle to capture highly variable and structured sequential data (e.g., multimodal behaviors) \cite{Zhan19}. 
Recent work in sequential generative models addressed this issue by injecting stochastic latent variables into the model and optimization using amortized variational inference to learn the latent variables (e.g., \cite{Chung15,Fraccaro16,Goyal17}). % 
%
Among these methods, a variational RNN (VRNN) \cite{Chung15} has been widely used in base models for multiagent trajectories \cite{Yeh19,Zhan19} with unknown governing equations.
% Note that our VRNN-based approach is compatible with other sequential generative models.
A VRNN is essentially a variational autoencoder (VAE) conditioned on the hidden state of an RNN and is trained by maximizing the (sequential) evidence lower-bound (ELBO):
}
\textcolor{black}{
\eq{
\mathcal{L}_{vrnn} =  &  \mathbb{E}_{q_{\phi}(z_{\leq T} \mid x_{\leq T})}\Bigg[ \sum_{t=1}^T \log p_{\theta}(x_t \mid z_{\leq T}, x_{<t}) \nonumber \\
& - D_{KL} \Big( q_{\phi}(z_t \mid x_{\leq T}, z_{<t}) || p_{\theta}(z_t \mid x_{<t}, z_{<t}) \Big) \Bigg], \label{eq:vrnn_elbo} 
}}
\textcolor{black}{where $z_t$ is a stochastic latent variable of VAE, and
$p_{\theta}(x_t \mid z_{\leq t}, x_{<t})$, $q_{\phi}(z_t \mid x_{\leq t}, z_{<t})$, and $p_{\theta}(z_t \mid x_{<t}, z_{<t})$ are generative model, the approximate posterior or inference model, and the prior model, respectively.
The first term is the reconstruction term.
The second term is the Kullback-Leibler (KL) divergence between the approximate posterior and the prior.}
% For the details, see Appendix \ref{app:vrnn}).

\textcolor{black}{
\noindent{\bf{GNN.}}
We then overview a graph neural network (GNN) based on \cite{Kipf18}.
Let $v_k$ be a feature vector for each node $k$ of $K$ agents. 
Next, a feature vector for each edge $e_{(k,j)}$ is computed based on the nodes to which it is connected. 
The edge feature vectors are sent as ``messages'' to each of the connected nodes to compute their new output state $o_k$.
Formally, a single round of message passing operations of a graph net is characterized below:
}
\textcolor{black}{
\eq{
v\rightarrow e: e_{(k,j)} &= ~f_e([v_k,v_j]), \\
e\rightarrow v: ~~~~~~o_i &= ~f_v\left(\sum_{j\in N(k)} e_{(k,j)}\right),
\label{eq:message} 
}}
\textcolor{black}{where $N(k)$ is the set of neighbors of node $k$ and $f_e$ and $f_v$ are neural networks.
In summary, a GNN takes in feature vectors $v_{1:K}$ and outputs a vector for each node $o_{1:K}$, i.e., $o_{1:K} = {\rm{GNN}}(v_{1:K})$. 
% , and an adjacency matrix $A$,
% Here, we consider fully connected graphs only, and thus do not use $A$ for simplicity. 
The operations of the GNN satisfy the permutation equivariance property as the edge construction is symmetric between pairs of nodes and the summation operator ignores the ordering of the edges \cite{zaheer2017deep}.
}

\vspace{10pt}
\section{Proposed Method}
\label{sec:proposed}
Here, we describe our \model method for ITE estimation in multiagent observational data. 
The overall framework is illustrated in Fig. \ref{fig:model}A.
% and Algorithm 1.
We aim to combine predictions of outcome and covariates using data-driven and theory-based approaches, while balancing the representations of treated and control groups to reduce the confounding bias.
To this end, we first introduce the representation learning of hidden confounders with balancing by mapping the current multiagent
observational data and historical information into the representation space.
Next, we describe the prediction methods of the time-varying covariates, a potential outcome, and the treatment using the learned representations.
Finally, we describe the loss function.

\subsection{Representation learning of confounders}
\label{ssec:representation}
\textcolor{black}{Here, as a main approach in Fig. \ref{fig:model}A, we extend} a GVRNN \cite{Yeh19} for local multiagent locations $x^l_t$  (i.e., specific for each agent) with theory-based computation.
\textcolor{black}{As its variant (e.g., used for the ablation study)}, a pure data-driven model combining GVRNN and VRNN \cite{Chung15} for global variables $x^g_t$ (i.e., common for all agents) is also considered. 
\textcolor{black}{Since the global variables do not usually have the graph structure, VRNN without GNN is suitable}. 
Here we describe the representation learning of hidden confounders.
% (note that Fig. \ref{fig:model} includes the theory-based computation instead of VRNN)
% We then explain the combined representation learning  VRNN, GVRNN, .

\noindent{\bf{GVRNN.}}
We first describe Graph Variational RNN (GVRNN) \cite{Yeh19} to obtain the representation from multiagent locations, which models the interactions between them at each step using GNNs. 
Let $x^l_{\leq T} = \{ x^l_1, \dots, x^l_T \}$ denote a sequence of covariates (here we consider multiagent locations). 
In this paper, GVRNN's update equations are as follows:
%
\eq{
p_{\theta}(z^l_t | x^l_{\leq t}, z^l_{<t}) & = \prod_k \mathcal{N}(z^l_{t,k}|\mu^{\rm{pri}}_{t,k},(\sigma^{\rm{pri}}_{t,k})^2), \\
q_{\phi}(z^l_t | x^l_{\leq t+1}, z^l_{<t}) & = \prod_k \mathcal{N}(z^l_{t,k}|\mu^{\rm{enc}}_{t,k},(\sigma^{\rm{enc}}_{t,k})^2),\\
p_{\theta}(x^l_{t+1} | z^l_{\leq t}, x^l_{\leq t}) & = \prod_k \mathcal{N}(z^l_{t,k}|\mu^{\rm{dec}}_{t,k},(\sigma^{\rm{dec}}_{t,k})^2),\\
h^l_{t+1,k} & = f^l_{rnn}(x^l_{t+1,k}, z^l_{t,k}, h^l_{t,k}),
\label{eq:gvrnn_state}
}
%
where $h^l_t$ and $z^l_t$ are deterministic and stochastic latent variables, $\mathcal{N}(\cdot|\mu,\sigma^2)$ denotes a multivariate normal distribution with mean $\mu$ and covariance matrix diag($\sigma^2$), and
%
\eq{
[\mu^{\rm{pri}}_{t,1:K},\sigma^{\rm{pri}}_{t,1:K}] & = {\rm{GNN_{pri}}}(h^l_{t,1:K}),\\
[\mu^{\rm{enc}}_{t,1:K},\sigma^{\rm{enc}}_{t,1:K}] & = {\rm{GNN_{enc}}}([x^l_{t+1,1:K},h^l_{t,1:K}]), \\
[\mu^{\rm{dec}}_{t+1,1:K},\sigma^{\rm{dec}}_{t+1,1:K}] & = {\rm{GNN_{dec}}}([z^l_{t,1:K},h^l_{t,1:K}]).
}
%
The prior network ${\rm{GNN_{pri}}}$, encoder ${\rm{GNN_{enc}}}$, and decoder ${\rm{GNN_{dec}}}$ are GNNs with learnable parameters $\phi$ and $\theta$.
\textcolor{black}{The relationship among them is illustrated in Fig. \ref{fig:model}B.}
Here we used the mean value $\mu^{\rm{dec}}_{t+1,1:K}$ as input variables $\hat{x}^{l'}_{t+1}$ in the following theory-based computation.
% Note that, while the RNN state for agent $k$ only depends directly on quantities related to $k$, these quantities do depend indirectly on the other agents through the GNNs. 
%In practice, the model generates the change in coordinates at each step $\Delta x_{t,k}$, and then computes $x_{t,k} = x_{t-1,k} + \Delta x_{t,k}$ \cite{Yeh19} (i.e., skip connection).
GVRNN is trained by maximizing the sequential ELBO in a similar way to VRNN as described in Eq. (\ref{eq:vrnn_elbo}), which is denoted as $\mathcal{L}_{gvrnn}$.

\noindent{\bf{Combined representation learning.}}
To construct a fully data-driven model (instead of the theory-based computation in Fig. \ref{fig:model}A), we propose a hierarchical GVRNN combining GVRNN for multiagent locations and VRNN for global inputs to learn the representation of hidden confounders.
In summary, each agent’s trajectory and other global information are processed through gated recurrent units (GRUs). 
The GRU parameters for the agent's trajectory are shared but keep its own individual recurrent state. 
At each time stamp, the model takes the current covariates and treatment assignments as input for learning representations of hidden confounders $z_t = [z^l_t,z^g_t]$ via GRUs, GNN, and MLP encoders. 

\subsection{Prediction with learned representation}
Our methods predict time-varying covariates, potential outcomes, and treatment for balancing by combining data-driven and theory-based approaches.
Here our contributions are to propose the theory-based computation of global variables and the prediction of time-varying covariates.
In this subsection, we describe the theory-based computation and the prediction of time-varying covariates, potential outcomes, and treatment for balancing.

\noindent{\bf{Theory-based computation.}} 
Here we assume that we do not have simulators including governing equations (used e.g., in  \cite{Kipf18,takeishi2021physics}) of multiagent systems such as team sports.  
In such a situation, we can utilize theory or prior knowledge of the domain using two approaches. 
One is to partially incorporate rule-based models into the data-driven model such as a mathematical relationship \cite{Fujii20policy} (e.g., between position and velocity), the biological constraints (e.g., turn angle in Sec. \ref{ssec:synthetic}), and critical behaviors (e.g., collision in Sec. \ref{ssec:synthetic}, ball movement in the air, and defending against the shot in Sec. \ref{ssec:basketball}).
Another is to compute auxiliary features for potential outcome predictions such as global variables (e.g., mean angular momentum in Sec. \ref{ssec:synthetic} or specific inter-agent distances in Secs. \ref{ssec:synthetic} and \ref{ssec:basketball}).
In our problem, we need to predict potential outcomes and covariates using predicted local variables from the data-driven model in Sec. \ref{ssec:representation} as shown in Fig. \ref{fig:model}A.
That is, 
\eq{\hat{y}_{t+1} = f^y_{theory}([z^l_t,x^g_t,a_t])\label{eq:theoryy}\\
\hat{x}_{t+1} = f^x_{theory}(\hat{x}^{l'}_{t+1}),\label{eq:theoryx}
}
% \mu^{\rm{dec}}_{t,1:K}
where $\hat{x}_{t+1} = [\hat{x}^l_{t+1},\hat{x}^g_{t+1}]$ and $f^y_{theory}$ and $f^x_{theory}$ are theory-based functions utilizing domain knowledge. 
For the details, see Sec. \ref{sec:experiments}.

\noindent{\bf{Prediction of time-varying covariates.}} 
Another contribution of this paper is to propose methods to predict time-varying covariates, which can confirm under what circumstances the intervention is effective.
The proposed \model infers time-varying covariates via GNN decoders and the theory-based computation as illustrated in Fig. \ref{fig:model}A.
For a fully data-driven model, we use the VRNN decoder for global variables instead of the theory-based computation.
We minimize the factual covariate loss function as follows:
\begin{equation}
\label{eq:covariate_loss}
    \mathcal{L}_x=\frac{1}{nT}\sum^{n}\sum_{t=1}^{T}(\hat{x}_{t}-x_{t})^{2}.
\end{equation}

\noindent{\bf{Prediction of a potential outcome.}} 
Next, we describe a potential outcome prediction network to estimate the outcome $\hat{y}_{t+1}$ as described in Eq. (\ref{eq:theoryy}).
Depending on the problem of various data domains, MLP was placed before or after the theory-based function $f^y_{theory}$ to properly model the potential outcome (for details, see Section \ref{sec:experiments}).
For a fully data-driven model, we use MLP instead of the theory-based computation.
We minimize the factual loss function as follows:
\begin{equation}
\label{eq:outcome_loss}
    \mathcal{L}_y=\frac{1}{nT}\sum^{n}\sum_{t=1}^{T}(\hat{y}_{t+1}-y_{t+1})^{2}.
\end{equation}

\noindent{\bf{Treatment prediction and balancing.}} 
We also predict the treatment assignments $\hat{a}_t$ at each time stamp. The predicted treatments $\hat{a}_t$ are obtained through a fully-connected layer with a sigmoid function as the last layer.
That is, $\hat{a}_t$ is the probability of receiving treatment based on the confounders at time $t$, which can be typically referred to as propensity score \cite{rosenbaum1983central} $\hat{a}_t=p(a_{t}=1|z_{t})$.

Since we consider the binary treatment in this paper, we use a cross-entropy loss for the treatment prediction as follows:
%\begin{small}
\begin{equation}
\label{eq:treatment-pred-loss}
    \mathcal{L}_a=-\frac{1}{nT}\sum^{n}\sum_{t=1}^{T}(a_t\log{\hat{a}_t}+(1-a_t)\log{(1-\hat{a}_t)}).
\end{equation}
%\end{small}

One critical consideration is balancing the representations of treated and control groups, which helps reduce the confounding bias \cite{bica2020estimating,liu2020estimating,ma2021deconfounding} and minimize the upper bound of the outcome inference error \cite{shalit2017estimating}. 
In this paper, we incorporate the
adversarial learning-based balancing method using a gradient reversal layer \cite{ganin2016domain} according to the related studies \cite{bica2020estimating,ma2021deconfounding}.
\textcolor{black}{Figure \ref{fig:GRL} illustrates our adversarial training process.
Let $\rm{GNN_{enc}(\cdot ; \theta_g)}$ be the GNN encoder with parameters $\theta_g$.
Also, let $\rm{MLP_a(\cdot ; \theta_a)}$ and $f^y_{theory}(\cdot ; \theta_y)$ be the MLP for treatment and outcome prediction with parameters $\theta_a$ and $\theta_y$ (if $f^y_{theory}$ is a neural network). $\gamma$ and $\lambda$ are the hyperparameters (or weights) of loss functions as described below. 
}
We add the gradient reversal layer before the treatment classifier \textcolor{black}{$\rm{MLP_a(\cdot; \theta_a)}$} to ensure the confounder representation distribution of the treated and that of the controlled are similar at the group level \cite{bica2020estimating,ma2021deconfounding}. 
The gradient reversal layer does not change the input during the forward-propagation, but in the back-propagation, reversing the gradient by multiplying by a negative scalar \textcolor{black}{(i.e., $-\lambda \frac{\partial \mathcal{L}_a}{\partial \theta_g}$ in Fig. \ref{fig:GRL}}). 
Although the model will minimize the loss of the treatment prediction, the adversarial learning process will contribute to the distribution balancing during the prediction of the potential outcome and the covariates. 

\begin{figure}[h]
\centering
\includegraphics[width=0.5\textwidth]{fig/Fig_GRL.png}
\caption[]{\textcolor{black}{Illustration of gradient reversal layer (GRL). The diagram is a part of Fig. \ref{fig:model} ignoring covariate prediction and includes backward passes. $\theta_g$, $\theta_a$, and $\theta_y$ are the neural network parameters of the GNN encoder, MLP for treatment and outcome prediction (if $f^y_{theory}$ is a neural network). 
GRL multiplies the gradient by a certain negative constant during the backpropagation-based training. }
}
%
\label{fig:GRL}
\end{figure}

\if0
\noindent{\bf{Outcome prediction constraint based on covariate prediction.}} 
Finally, we model an outcome prediction constraint based on covariate prediction for reducing the discrepancies between both predictions.
In our problem, the predicted outcome at time $t+1$ (here, we denote $\tilde{y}_{x,t+1}$) can be computed from $\hat{x}_{t+1}$ but in our model, $\hat{y}_{t+1}$ is not generated from  $\hat{x}_{t+1}$. 
We then minimize a loss function for reducing the discrepancies as follows:
\begin{equation}
\label{eq:covariate_loss}
    \mathcal{L}_{yx}=\frac{1}{n(T-1)}\sum^{n}\sum_{t=1}^{T-1}(\hat{y}_{t+1}-\tilde{y}_{x,t+1})^{2}.
\end{equation}
\fi

\subsection{Loss function}
The loss function for our method is defined as
\begin{equation}
    \mathcal{L} = \mathcal{L}_y + \alpha \mathcal{L}_{gvrnn} + \gamma \mathcal{L}_x + \lambda \mathcal{L}_a %+ \ \mathcal{L}_{yx}, 
    \label{eq:loss}
\end{equation}
where $\mathcal{L}_y$ is the factual outcome prediction loss, $\mathcal{L}_{gvrnn}$ is the ELBO in GVRNN, $\mathcal{L}_x$ is the covariate prediction loss, and $\mathcal{L}_a$ is the treatment prediction loss. %, and $\mathcal{L}_{yx}$ is an outcome prediction constraint based on covariate prediction.
$\alpha$, $\gamma$, and $\lambda$ 
are hyperparameters to balance the loss function. 
In a pure data-driven model, we add $\mathcal{L}_{vrnn}$ to $\mathcal{L}$.
To prevent over-fitting, we select the best-performing model using the loss function on the validation set.
The sensitivity analysis in the hyperparameters is presented in Appendix \ref{app:sensitivity}.



\section{Related work}
\label{sec:related}
\noindent\textbf{Learning causal effects with time-varying data.}
% survey \cite{moraffah2021causal}
Pioneering work to estimate the effects of time-varying exposures has developed such as 
g-computation, structural nested models, and marginal structural models (e.g., \cite{robins1986new, robins1994correcting, robins2000marginal, robins2009estimation}) in statistics and epidemiology domains. 
To handle complex time-dependencies, Bayesian non-parametric approaches based on Gaussian processes \cite{xu2016bayesian,schulam2017reliable,soleimani2017treatment} and Dirichlet processes \cite{roy2017bayesian} have been proposed. 
Moreover, to model time-dependencies without strong assumptions on functional forms, RNN approaches have been intensively investigated such as recurrent marginal structural networks \cite{lim2018forecasting} and adversarial training to balance the historical confounders \cite{bica2020estimating}. 
Recently, to estimate the treatment effects with hidden confounders, several methods have been proposed by relaxing the assumption of strong ignorability \cite{bica2020time,liu2020estimating,ma2021deconfounding}.
Our approach is related to these works, but they did not consider the utilization of domain knowledge and the explicit prediction of time-varying covariates, which is necessary for the proper interpretation of the results.
% confirming under what situation the treatment is effective.

\noindent\textbf{Representation learning for treatment effect estimation.}
Estimation methods of balanced representation between treated and control groups in hidden space have been proposed in the static setting \cite{hill2011bayesian,wager2018estimation,alaa2018limits}.
In neural network approaches, methods with regularization for the balancing \cite{johansson2016learning,shalit2017estimating}, incorporating the local similarity among individuals \cite{yao2018representation},  generative adversarial nets (GAN) \cite{yoon2018ganite}, a multi-task learning \cite{shi2019adapting}, \textcolor{black}{probabilistic 
 modeling \cite{grecov2022probabilistic,abroshan2022conservative}, and optimal transport framework \cite{li2021causal}} have been proposed.  
In the dynamic setting, \cite{bica2020estimating,ma2021deconfounding} adopted adversarial training techniques with a gradient reversal layer, \cite{bica2020time} proposed multi-task learning to build a factor model of the cause distribution, and \cite{liu2020estimating} used inverse probability of treatment weighting. 
\textcolor{black}{To model multiagent or networked systems, GNN-based approaches have been intensively used in prediction problems (e.g., \cite{wu2019graph, zhang2020spatio}).} 
To obtain representations from the networked covariates, \cite{ma2021deconfounding} incorporated graph convolutional networks \textcolor{black}{and \cite{takeuchi2023causal} incorporated hierarchical spatial graph structure into causal estimation frameworks}
(other counterfactual learning on graphs was surveyed by \cite{guo2023counterfactual}). Recently, continuous-time causal inferences for dynamical systems \cite{seedat2022continuous,jiang2023cf} have been proposed.
We adopted the adversarial training to balance the representation and introduced GVRNN with prior (domain) knowledge for a more accurate long-term prediction of potential outcomes and covariates.

% adversalial \cite{ganin2016domain}

\noindent\textbf{Counterfactual prediction for multiagent data.}
In past research of real-world multiagent movements, counterfactual prediction of human movements has been studied.
\textcolor{black}{As control problems, forward models such as in robotics, can be considered. However, our problem is inherently a backward problem to model the multiagent behaviors from data, then here we focus on such backward approaches.}
Counterfactual trajectory prediction methods have been proposed in pedestrians \cite{chen2021human}, animals \cite{fujii2021learning}, and team sports \cite{Yeh19,Fujii20policy,nakahara2022estimating,teranishi2022evaluation}.
Outside the context of causal inference based on the Rubin causal model, the causality of the physics-based system such as using neural rational inference \cite{Kipf18} and learning physics in a game engine with randomized stability \cite{lerer2016learning} have been considered (sometimes called causal reasoning).
\textcolor{black}{In the context of reinforcement learning, the counterfactual action prediction and evaluation can be considered in team sports \cite{Liu2018,Liu2020,rahimian2022beyond,nakahara2023action}.}
In traffic environments, 
methods to understand the reasons behind the maneuvers of other vehicles and pedestrians for escaping accidents \cite{ramanishka2018toward, you2020traffic,mcduff2021causalcity} have been proposed.
In the static setting of causal inference, \cite{takeuchi2021grab} proposed a spatial convolutional counterfactual regression to estimate the effects of crowd movement guidance. % takeuchi 2023?
In team sports, propensity score matching was used to investigate the causal effect of some plays or timeouts in many sports \cite{yam2019lost,toumi2019grapes,gibbs2020causal,nakahara2022pitching}.
% going for the touchdown in American football \cite{yam2019lost}, % the NFL, going for it on fourth down
% clearing the puck in ice hockey \cite{toumi2019grapes}, % NHL
% and the effectiveness of timeouts in basketball  \cite{gibbs2020causal}.
% at stopping an opposing run
% crossing the ball in soccer \cite{wu2021contextual}.
% baseball pitching \cite{nakahara2022pitching}.
In the dynamic setting, \cite{vock2018estimating} applied a g-computation method to examine the effect of a specific pitch in baseball. %taking a pitch during a 3-0 count in MLB. 
We firstly propose a framework for estimating ITE in multiagent motions. % for the first time. %  causal inference 

\section{Experiments}
\label{sec:experiments}
The purpose of our experiments is to validate the proposed methods for application to real-world multiagent trajectories, which usually have no ground truth of the interaction rules. 
Hence, for verification of our methods, we first compared their performances to infer the ITE to those in various baselines using two synthetic datasets with ground truth: the CARLA \cite{dosovitskiy2017carla} autonomous driving simulator and a biological multiagent model called Boid \cite{Couzin02}.
In particular, for real-world applications, we need long-term counterfactual predictions.
We mainly verified the estimation of ITE using the CARLA dataset (relatively fewer interactions) and that of the best intervention timing using the Boid dataset (relatively more and complex interactions) because the intervention timing is sensitive.
Finally, we examined the applicability of our methods to real-world data using the basketball (NBA) dataset.
In common, we separated the time $T$ into a burn-in period $1,\ldots,T_b$ and a prediction period $T_p = T_b+1,\ldots,T$.
The intervention point during the intervention period is denoted as $T_i$ (a similar interval to $T_p$, but we set it differently in each problem). 
%  (here binary treatments) 
The hyperparameters of the models were determined by validation datasets in each experiment (for the details, see Appendices \ref{app:boid}, \ref{app:carla}, and \ref{app:nba}). 

Here, we commonly compared our methods to four baseline methods: 
a simple RNN baseline using GRU \cite{Cho14},
deep sequential weighting (DSW) \cite{liu2020estimating} as a baseline considering hidden confounders, graph counterfactual recurrent network (GCRN: for clarity, we change the name) \cite{ma2021deconfounding} 
%(replacing VRNN with RNN and not predicting time-varying covariates)
, and the variant to predict covariates (GCRN$+$X).
These baselines were modified to our setting (e.g., multiagent and time-varying outcome: for details, see Appendix \ref{app:common}). 
For verification, since the prediction of the covariates is required, the most appropriate baseline is GCRN$+$X, which is compared via visualization because DSW and GCRN cannot predict the covariates.  
We also validated our approach with three variants: GV-CRN (without the theory-based computation, i.e., without global covariates), TV-CRN (removing GNN), and TG-CRN (replacing VRNN with RNN).
% First, we evaluated the effectiveness of GNN and VRNN using the first and second variants. 
\textcolor{black}{To perform fair comparisons among the models, we trained all models with 20 epochs (the learning curves are shown in Appendix \ref{app:converge}).}
Other common training details 
% and computational resources
% , and the amount of computation 
are described in Appendix \ref{app:common} and the codes are provided at \url{https://github.com/keisuke198619/TGV-CRN}.


\subsection{Synthetic datasets}
\label{ssec:synthetic}
\vspace{-0pt}
To verify our method, we compared the performances to infer the causal effect with those in various baselines using two synthetic datasets with ground truth.
\textcolor{black}{We used two types of simulations: autonomous vehicle and biological agent (Boid) simulations. All simulations are rule-based rather than learning-based as described below.}
As for performance metrics, we first adopted commonly used potential outcome and covariate prediction errors.
Potential outcome prediction errors were computed as an absolute error of all simulated outcomes: $L_{outcome}^{(a_t,t,i,t')} = |\hat{y}_{a_t,t}^{(i,t')} - y_{a_t,t}^{(i,t')} |$, where a treatment $a_t = \{0,1\}$, time $t = [T_b,T+1]$, simulated temporal variation of the treatment timing $t' \in T_i$, and the burn-in period in RNN $T_b$.
Covariate prediction errors were computed as $L_2$ prediction error of all simulated covariates: $L_{covariates}^{(a_t,t,i,t')} = \|\hat{x}_{a_t,t}^{(i,t')} - x_{a_t,t}^{(i,t')} \|_2$.
% $L_2^{n,m} = (1/(2T\cdot K))\sum^T _{t=1}\sum^K_{k=1}\|\hat{a}^n_{m,t,k} - a_{m,t,k} \|_2$
We took the average over all variables for evaluation.

We also adopted two widely-used evaluation metrics in causal inference: rooted precision in the estimation of heterogeneous effect (PEHE) \cite{hill2011bayesian} 
$\sqrt{\epsilon^t_{PEHE}}=\sqrt{\frac{1}{n}\sum_{i\in n}(\hat{\tau}_t^{(i)}-\tau_t^{(i)})^2}$
and mean absolute error of the average treatment effect (ATE) \cite{willmott2005advantages} to measure the quality of the estimated individual treatment effects at different time stamps: $\epsilon^t_{ATE}=|\frac{1}{n}\sum_{i\in n}\hat{\tau}_t^{(i)}-\frac{1}{n}\sum_{i\in n}\tau_t^{(i)}|$.
% as follows.
% \eq{\sqrt{\epsilon^t_{PEHE}}=\sqrt{\frac{1}{n}\sum_{i\in n}(\hat{1}_t^{(i)}-1_t^{(i)})^2},\label{eq:pehe}}
% \eq{\epsilon^t_{ATE}=|\frac{1}{n}\sum_{i\in n}\hat{1}_t^{(i)}-\frac{1}{n}\sum_{i\in n}1_t^{(i)}|.\label{eq:ate}}
We took the average over all time stamps for evaluation.

\textcolor{black}{Regarding interpretability, we emphasize that there has been no previous work to visualize or interpret future covariates (i.e., multiagent trajectory), which can confirm under what circumstances the intervention is effective. 
We then illustrated trajectory prediction results in each domain.
We compared the customized baselines to predict future covariates to demonstrate the superiority of our approach, which are also quantitatively shown as the covariate losses $L_{covariates}$. }

\noindent {\bf{Autonomous vehicle simulation.}}
We first validated the performances of our methods using an autonomous vehicle simulation using the CARLA simulator \cite{dosovitskiy2017carla} (ver. 0.9.8). 
The CARLA environment contains dynamic obstacles (e.g., pedestrians and cars) that interact with the ego car. 
To generate data, we simulated autonomous driving at various towns, starting points, and obstacle positions, which were subsampled at 2 Hz (see also Appendix \ref{app:carla}).
Since we did not use the future path as inputs for generality, we randomly split 7 towns data into 904 training, 129 validation, and 259 test scenarios. % total 1292
Here we used two types of autonomous vehicles: partial observation with Autoware \cite{kato2015open} (ver. 1.12.0) and full observation types with a pre-installed CARLA simulator. 
The partial observation model often stops when dangerous situations arise for safety. 
The full observation model intervenes to accelerate the stopped vehicle after the safety confirmation. 
We set $T=60, T_b=40$, and $T_i=T_b+1,\ldots,T_b+10$. %, and $1 = 1$.
For counterfactual data, since dangerous situations (i.e., requiring intervention) are limited, we only used with or without intervention at the same timing (not various timings). % where intervention is required 

We predict the safe driving distance of the ego car from the starting point as an outcome while driving without any collision with obstacles. 
That is, the treatment effect of the full observation model is defined as the difference between safe driving distances with and without interventions.
As the local covariates $x^l_t$, we used the position, velocity, pose, and size information of the ego car and obstacles in the 2D map.
The maximum number of dynamic obstacles was $120$, but since the obstacles related to the ego car were limited and the number changes over time, we used the nearest $10$ obstacles' information as the covariates. 
As the global covariates $x^g_t$, we used the current driving distance $x_t^{dist}$ and (binary) collision $x_t^{coll}$ information of the ego car.
The theory-based function $f^x_{theory}$ mathematically computed $\hat{x}_{t+1}$ using the predicted velocity and pose information, and added the learned positive velocity in intervention.
$f^y_{theory}$ computed $\hat{y}_{t+1}$ using the predicted $\hat{y}'_{t+1}$ as the output of the MLP such that $\hat{y}_{t+1} = (1-x_t^{coll})\hat{y}'_{t+1}$ based on the definition of the safe driving distance.

Quantitative verification results are shown in Table \ref{tab:carla}.
%Here we show the average prediction errors of the outcome and covariates over all samples and time stamps, and the average $\sqrt{\epsilon^t_{PEHE}}$ and $\epsilon^t_{ATE}$ over all time stamps.
Our full model and its ablated variants show lower prediction errors of the outcome and covariates, and $\sqrt{\epsilon^t_{PEHE}}$, and $\epsilon^t_{ATE}$ than other baselines. 
% For the $\sqrt{\epsilon^t_{PEHE}}$, our model, its variants, and GCRN \cite{ma2021deconfounding} in our setting show competitive performances, which were better than those of DSW \cite{liu2020estimating} without GNN and the simple RNN. 
Totally, our full model and the variant without amortized variational inference (TG-CRN) show the best performances, suggesting that the necessity of complex modeling for long-term prediction using this dataset would be smaller than that using the following datasets. 
In our four models, we found that on this dataset the combination of theory-based computation and GNN worked well in the covariate prediction, but did not in the outcome prediction.  
Example results of our method are shown in Fig. \ref{fig:carla}\textcolor{black}{, which can interpret the results}.
Our method (\model) shows better counterfactual prediction with intervention than the baseline (GCRN$+$X).
On average, the outcome values with intervention compared to the absence of the intervention (i.e., ITE) were $0.005 \pm 0.000$ (km) in our full model, $-0.021 \pm 0.000$ in the baseline (GCRN+X), and $0.013 \pm 0.001$ in the ground truth.
As Fig. \ref{fig:carla} shows, the baseline did not model the intervention effects.
\textcolor{black}{We expect that this method can be applied to the effect of human intervention in Level 3 autonomous vehicle simulations, which need human intervention. This approach can examine the effect of autonomous control in Level 4 or 5 autonomous vehicle simulations (without human intervention) in a case when human intervention is necessary in Level 3.}

% based on the ground truth without intervention.
\newcommand{\md}[2]{\multicolumn{#1}{c|}{#2}}
\newcommand{\me}[2]{\multicolumn{#1}{c}{#2}}
\begin{table}[ht!]
\centering
\scalebox{0.8}{
\begin{tabular}{l|cccc}%|
\Xhline{3\arrayrulewidth} %\hline
% & \me{4}{carla simulation dataset}  \\ 
& \me{1}{$L_{Outcome}$ } & \me{1}{$\sqrt{\epsilon^t_{PEHE}}$} &  \me{1}{$\epsilon^t_{ATE}$} & \me{1}{$L_{Covariates}$} \\
\hline
RNN & 3.330 $\pm$ 0.286 &  0.159 $\pm$ 0.025 & 0.129 $\pm$ 0.023 & 0.479 $\pm$ 0.0011
%\\CRN \cite{bica2020estimating} &   &  & 
\\DSW \cite{liu2020estimating}& 0.161 $\pm$ 0.013 & 0.030 $\pm$ 0.003 & 0.019 $\pm$ 0.003  & ---
\\GCRN \cite{ma2021deconfounding} &  0.094 $\pm$ 0.009 & 0.024 $\pm$ 0.002 & 0.014 $\pm$ 0.002 & ---
\\GCRN$+$X &  3.020 $\pm$ 0.496 &  0.049 $\pm$ 0.009 & 0.034 $\pm$ 0.007 & 2.072 $\pm$ 0.0098\\
\hline
GV-CRN & \textbf{0.038} $\pm$ \textbf{0.003} &  0.022 $\pm$ 0.002 & 0.012 $\pm$ 0.001 & 0.372 $\pm$ 0.0010 
\\TG-CRN &  0.045 $\pm$ 0.004 &  \textbf{0.021} $\pm$ \textbf{0.002} & \textbf{0.011} $\pm$ \textbf{0.001} & \textbf{0.082} $\pm$ \textbf{0.0004}
\\TV-CRN & \textbf{0.032} $\pm$ \textbf{0.003} &  0.023 $\pm$ 0.002 & 0.014 $\pm$ 0.002 & 0.102 $\pm$ 0.0004 
\\ TGV-CRN (full) & 0.041 $\pm$ 0.003 &  \textbf{0.020} $\pm$ \textbf{0.002} & \textbf{0.008} $\pm$ \textbf{0.001} & \textbf{0.098} $\pm$ \textbf{0.0004}
\\
\Xhline{3\arrayrulewidth} % \hline
\end{tabular}
}
\vspace{3pt}
\caption{\label{tab:carla} Performance comparison on the carla dataset. \textcolor{black}{The upper and lower rows indicate the baselines and proposed methods, respectively.}}
% of  model simulation.}
%Standard error are evaluated across 400 test samples.}
\end{table}

\begin{figure}[h]
\centering
\includegraphics[width=0.45\textwidth]{fig/Fig_carla.png}
\caption[]{Example CARLA results using our method. (Top) Visualization of covariates and (\textcolor{black}{middle row and} bottom) outcome time series in (left) ground truth without intervention, (middle \textcolor{black}{column}) counterfactual intervention using our model, and (right) the baseline. 
The middle \textcolor{black}{row} subfigures are enlarged views of the bottom ones \textcolor{black}{from 20 s}.
An ego car (red square) and obstacles (black) are shown in the upper plots (see also Fig. \ref{fig:intervention}A) at the intervention time, which is the solid line in the lower plots.
% Red square is an ego car, red line is the velocity, filled square is the start point, and
The unfilled circle is the start of the long-term prediction (dashed line in the lower plots).
The ego-car moves from right to left and stops because of the obstacles.% Solid line in lower plots is the intervention time.
\textcolor{black}{The videos are given in the above GitHub page.}
}
\label{fig:carla}
\end{figure}


\begin{table*}[h!]
\centering
\scalebox{0.9}{
\begin{tabular}{l|ccccc|ccc}%|
\Xhline{3\arrayrulewidth} %\hline
& \md{5}{Boid simulation dataset}& \me{3}{Real-world NBA dataset}  \\ 
&\me{1}{Treatment timing} & \me{1}{$L_{Outcome}$ } &
\me{1}{$\sqrt{\epsilon^t_{PEHE}}$} &  \me{1}{$\epsilon^t_{ATE}$}  & \md{1}{$L_{Covariates}$} &   \me{1}{$L_{Outcome}$ } & \me{1}{$\hat{\tau}^{CFP}$} & \me{1}{$L_{Covariates}$} \\ 
\hline
RNN & 1.815 $\pm$ 0.068 & 0.668 $\pm$ 0.033 &  0.443 $\pm$ 0.054 & 0.085 $\pm$ 0.016 & 0.168 $\pm$ 0.0002 & 0.291 $\pm$ 0.003 &  0.170 $\pm$ 0.0065&  0.969 $\pm$ 0.0024
%\\CRN \cite{bica2020estimating} &   &  & 
\\DSW \cite{liu2020estimating}
& 2.353 $\pm$ 0.080 & \textbf{0.586} $\pm$ \textbf{0.029} & \textbf{0.440} $\pm$ \textbf{0.054} & \textbf{0.082} $\pm$ \textbf{0.015} & --- & 0.197 $\pm$ 0.003&  0.166 $\pm$ 0.0064 & ---
\\GCRN \cite{ma2021deconfounding} 
& 2.290 $\pm$ 0.081 & \textbf{0.587} $\pm$ \textbf{0.028} & 0.443 $\pm$ 0.055 & 0.088 $\pm$ 0.016 & --- & \textbf{0.170 }$\pm$ \textbf{0.003}&  0.149 $\pm$ 0.0062  & ---
\\GCRN$+$X & 2.290 $\pm$ 0.081 & 0.727 $\pm$ 0.022 &  \textbf{0.440} $\pm$ \textbf{0.054} & \textbf{0.080} $\pm$ \textbf{0.014} & 0.162 $\pm$ 0.0001 & 0.431 $\pm$ 0.004 &  0.154 $\pm$ 0.0063&  1.226 $\pm$ 0.0032\\
\hline
GV-CRN &  1.900 $\pm$ 0.068 & 0.674 $\pm$ 0.028 &  0.536 $\pm$ 0.045 & 0.090 $\pm$ 0.014 & 0.329 $\pm$ 0.0003 & 0.475 $\pm$ 0.005 &  0.151 $\pm$ 0.0060&  1.173 $\pm$ 0.0028
\\TG-CRN & \textbf{1.750} $\pm$ \textbf{0.062} & 0.694 $\pm$ 0.041 &  0.501 $\pm$ 0.067 & 0.125 $\pm$ 0.021 & \textbf{0.086} $\pm$ \textbf{0.0002} & 0.347 $\pm$ 0.005 &  0.205 $\pm$ 0.0072&  \textbf{0.811} $\pm$ \textbf{0.0022}
\\TV-CRN & \textbf{1.692} $\pm$ \textbf{0.064} & 0.981 $\pm$ 0.044 &  0.483 $\pm$ 0.052 & 0.106 $\pm$ 0.018 & 0.090 $\pm$ 0.0002  &  \textbf{0.165} $\pm$ \textbf{0.002} &  0.192 $\pm$ 0.0068&  0.852 $\pm$ 0.0023
\\ TGV-CRN (full) & 1.853 $\pm$ 0.062 & 0.690 $\pm$ 0.032 &  0.537 $\pm$ 0.055 & 0.101 $\pm$ 0.016 & \textbf{0.085} $\pm$ \textbf{0.0002} & 0.231 $\pm$ 0.004 &  0.261 $\pm$ 0.0083&  \textbf{0.825} $\pm$ \textbf{0.0023}
\\
\Xhline{3\arrayrulewidth} % \hline
\end{tabular}
}
\vspace{3pt}
\caption{\label{tab:boid} Performance comparison on the boid and NBA datasets.\textcolor{black}{The upper and lower rows indicate the baselines and proposed methods, respectively.}}
% of  model simulation.}
%Standard error are evaluated across 400 test samples.}
\end{table*}

%\vspace{-5pt}
\noindent {\bf{Biological multiagent simulation.}}
Here, we validated our methods on the Boid model, which contains movement trajectories of 20 agents.
% (for details, see Appendix \ref{app:boid}). 
The Boid model (originally, \cite{reynolds1987flocks}) is a rule-based model to generate generic simulated flocking agents and we used a unit-vector-based (rule-based) model \cite{Couzin02} (for details, see Appendix \ref{app:boid}).
In this paper, we intervene the agents' recognition to generate torus (circle) behaviors from the swarm (random) behaviors.
The outcome is defined as the mean angular momentum of individuals about the center of the group (assuming the mass of each agent $m=1$). % (see also Appendix \ref{app:boid})
That is, the treatment effect of the change in the recognition is estimated as the difference in the future mean angular momentum between simulations with and without the interventions.

In this model, $20$ agents are described by a 2-D vector with a $1$ m/s constant velocity in a 15 $\times$ 15 m boundary square.
At each time stamp, a member will change direction according to the positions of all other members based on three zones.
The first is the repulsion zone with radius $r_r = 0.5$ m, in which individuals within each other’s repulsion zone try to avoid each other by swimming in opposite directions. 
The second is the orientation zone, in which individuals try to move in the same direction; 
here we set radius $r_o = 1$ to generate swarming behaviors before the intervention. 
To generate torus behaviors, we change $r_o = 4$, which is the intervention in this study.
The third is the attractive zone (radius $r_a = 7.5$ m), in which agents move towards each other and tend to cluster.
%, while any agents beyond that radius have no influence. 

To simulate the treatment assignments, we generate factual $20800$ samples ($20000$ training, $400$ validation, and $400$ test datasets).
We set $T=14$ and $T_b = 9$ and randomly pick the intervention point during the intervention period $T_i = 9,\ldots,13$. 
The outcome is defined as the mean angular momentum among individuals at time $t+1$. % (i.e., $\tau=1$).
We also created a counterfactual dataset only in the test dataset.
Here we created all combinations of treatment points during the intervention period $T_i$. 
As the local covariates $x^l_t$, we used position, velocity, and directional change of all agents.
As the global covariates $x^g_t$, we used the current mean angular momentum.
The theory-based function $f^x_{theory}$ mathematically computed $\hat{x}_{t+1}$ using the direction change $d$ at the next step, maximum turn angle $\beta$ as body constraints (for $d$ and $\beta$, see also Appendix \ref{app:boid}), and attraction rule when agents are far from the center of the group.
In addition, we added the orientation rule when agents are in the orientation zone and not in the repulsion zone in $f^x_{theory}$. 
$f^y_{theory}$ was replaced with a MLP such that $\hat{y}_{t+1} = {\rm{MLP_y}}([z^l_t,x^g_t,a_t])$.


% That is, the number of counterfactual samples is $500 \times 6$, where $6$ indicates all $5$ time points in $T_i$ and no treatment point except for the factual treatment point.  

The results are shown in Table \ref{tab:boid} left. 
In addition to the four indices in the CARLA experiment, we investigated the estimation error of the best intervention timing $|\argmax_{t'}(\hat{y}^{(i,t')}_{T+1})-\argmax_{t'}(y^{(i,t')}_{T+1})|$.
The results indicate that our model and its variants with theory-based computation show better prediction performances in covariates and the best intervention timing than all of the baselines.
%, and better than the most appropriate baseline (GCFN+$X$) in outcome prediction and $\epsilon^t_{ATE}$.
However, the outcome prediction errors in our models were worse than the causal inference baselines (DSW, GCRN, and GCFN+$X$), which may lead to degraded performances in  $\sqrt{\epsilon^t_{PEHE}}$ and $\epsilon^t_{ATE}$ of our models than the baselines. 
In our four models, all combination of core three components (T, G, and V) did not work well and on this dataset, TG-CRN without VRNN is the best performing model. 
On average, the ITE values were $0.160 \pm 0.010$ in our best model (TG-CRN), $0.035 \pm 0.001$ in the baseline (GCRN+X), and $0.091 \pm 0.026$ in the ground truth.
From this viewpoint (both average and variance), our best model was closer to the ground truth than the baseline. 
Example \textcolor{black}{interpretable} results of our method are shown in Fig. \ref{fig:boid}.
% , which indicates better and worse performances in our model. 
Our best model (TG-CRN) shows better counterfactual covariate prediction with intervention than the baseline (GCRN$+$X).
Moreover, the counterfactual prediction of the outcome in our model had variation among various intervention times, whereas that in the baseline did not.
% but the values were similar to the ground truth in average.
% the better counterfactual covariate prediction and the variation with various intervention timing 
We consider that these were important properties in our problem, and improved prediction of the potential outcome was left for future work. 
% but probably due to the difficulty for long-term outcome prediction, better
\textcolor{black}{We expect that our approach can estimate the effect of an experimenter’s interventions on multi-animal behaviors even if they did not perform the experiment in missing conditions. Our approach is expected to improve the efficiency of experimental procedures for observing desired movements.}

\begin{figure}[h]
\centering
\includegraphics[width=0.45\textwidth]{fig/Fig_boid.png}
\caption[]{Example Boid results of our method. The configurations are the same as Fig. \ref{fig:carla}. 
(Top) 20 boids with 6 different colors move in a rule-based manner (but we did not use this information for the prediction) and the filled circles are the starting point of the trajectories. 
(\textcolor{black}{Middle row and bottom) 
outcome time series in (left) ground truth without intervention, (middle column) counterfactual intervention using our model, and (right) the baseline. 
The middle row subfigures are enlarged views of the bottom ones from the 9th frame. }% The numbers are the start of the long-term prediction (trajectories). 
The ``a'' in the lower caption is the intervention times.
\textcolor{black}{For example, ``$a = 9$'' means the case of intervention at the 9th frame and ``None'' indicates no intervention.
}
\textcolor{black}{The videos are given in the above GitHub page.}
}
\label{fig:boid}
\end{figure}

\subsection{Real-world basketball dataset}
\label{ssec:basketball}
Finally, we examined the applicability of our methods to a real-world basketball dataset from the NBA.
Data acquisition was based on the contract between the league (NBA) and the company (STATS LLC.), not between the players and us. They are top-level players and then the data was not anonymized. The company was licensed to acquire this data, and it was guaranteed that the use of the data would not infringe on any rights of players or teams.
In this study, we used attack sequences from 630 games from the 2015/2016 NBA season (\url{ https://www.stats.com/data-science/}), which contained the trajectories of 10 players and the ball. 
We extracted 47,467 attacks (i.e., offensive plays) as samples, which were subsampled at 5 Hz.
We separated the dataset into 34,696 pre-training (458 games for training the following classifier of effective attack), 11460 training (154 games, 1/10 of that is used as validation), and 1,305 test samples (18 games) in chronological order.
%
Since scoring predictions are difficult in general (e.g., \cite{Fujii17,Fujii18}), 
we define the attack effectiveness as the outcome by predicting whether the attack is effective or not in the future.
This is because evaluating team movements based on scores alone may not provide a holistic view, due to factors such as the shooting skills of individual players.
In addition, we predict the attack effectiveness at the next time stamp using the pre-training samples and a logistic regression. % (i.e., $\tau=1$)   
The details are provided in Appendix \ref{app:nba}.
% including the rule-based definition of the effective attack 

We verified our methods using factual data and provided insights using counterfactual prediction.
In model verification using factual data, we set $T=95$ and $T_b = 85$.
For counterfactual predictions, we set $T=105$ and $T_b = 95$ (at the end of attacks, i.e., a shot or turnover) and predicted all combinations of the counterfactual timing during the intervention period $T_i = 95,\ldots,98$. 
As the local covariates $x^l_t$, we used the position and velocity of all agents including the ball.
As the global covariates $x^g_t$, we used the ball player's information and areas (see also Appendix \ref{app:nba}), distances from the nearest defender (about the ball player and other attackers), successful shot probabilities of all attackers, and game and shot clock.
The theory-based function $f^x_{theory}$ and $f^y_{theory}$ computed %$\hat{x}_^l{t+1}$ mathematically using the predicted velocity and
$\hat{x}^g_{t+1}$ and $\hat{y}_{t+1}$ in rule-based manners.
In addition, to perform long-term prediction, $f^x_{theory}$ also computed the ball and the defender nearest to the next ball player during a counterfactual pass in rule-based manners.

Our verification results using factual data are shown in Table \ref{tab:carla} right.
We examined the counterfactual pass effect $\hat{\tau}^{CFP} = \max_{t'\in T_i}(\max_{t=[T_b,T]}(\hat{y}^{(i,t')}_{T+1}))-y^{(i)}_{T_b+1}$ in addition to the outcome and covariate prediction errors. 
Our methods show better prediction performances in the factual covariates than the fully data-driven baselines. In the factual outcomes, our methods outperformed the most appropriate baseline (GCFN$+$X) and RNN, but only the method without GNN shows competitive performance with DSW and GCRN.  
One of the possible reasons may be the difficulty in modeling the relationship between the future outcome and current covariates, %  using GNN
whereas GNN worked well in only the covariate prediction as the previous work \cite{Yeh19}.
\textcolor{black}{As additional analysis, we indicate endpoint errors and distributions of long-term covariate prediction for basketball data in Appendix \ref{app:analysis_basket}. Again, the strength of our method is to model the covariates at the next timestep for interpretability of the model. Results in Appendix \ref{app:analysis_basket} %in Table \ref{tab:nba_end} 
that the tendency of the endpoint prediction error for all models is similar to the mean prediction error, in which our approach (TV-CRN and TGV-CRN) show the best performance. In addition, the distribution of the player velocity in the figure of Appendix \ref{app:analysis_basket} % . \ref{fig:histogram_nba} 
shows that our approach without theory-based computation had a wider distribution like ground truth than other models, and those with theory-based computation had less zero-velocity bins like ground truth than other models. Although our approach did not completely model the player’s velocity, we show that they were partially effective for covariate prediction using each component.}
% for long-term prediction in the outcome. 
In $\hat{\tau}^{CFP}$, since all models show positive values, it suggests that there may be a more promising shot opportunity by an extra pass in the shot situation. % of our models and baselines 
Figure \ref{fig:nba} shows realistic counterfactual prediction from our model 
\textcolor{black}{to demonstrate interpretable results.}
Our model completed the counterfactual pass and the nearest defender chase, but the baseline model failed and predicted unrealistic behaviors 
(e.g., the attackers moved toward the outside of the court).
\textcolor{black}{In practice, our approach can estimate the effect of the selection of passes in basketball shot scenarios. We expect that our approach can evaluate the decision-making skills of players in a competitive game.}
%(e.g., the defenders avoided the ball).
% the attacker \# 3 went outside of the court). 
% Although there was small variation, the counterfactual prediction of the potential outcome (effectiveness) in our model had variation among various intervention times (pass timing).
Again, we emphasize these important properties in our problem. 

\begin{figure}[h]
\centering
\includegraphics[width=0.45\textwidth]{fig/Fig_nba.png}
\caption[]{Example NBA results of our method. The configurations are the same as Fig. \ref{fig:boid}. 
(Top) Red and blue numbers, gray line, and orange circle and line indicate an attacker, a defender, players' historical trajectories, and the ball, respectively.
The positions of the numbers are at the end of the factual data (shot), which is shown as the break line in lower plots. 
In the CF intervention subplots, colored trajectories indicate counterfactual predictions.
The actual red player \#5 shot (left) but in the counterfactual prediction (middle and right \textcolor{black}{columns}), the player tried to pass to a teammate. In the middle top (our method), the player successfully passes to the teammate red \#3, but in the left top (baseline), the player’s pass failed. 
\textcolor{black}{(Bottom) 
outcome time series (attack effectiveness) are shown. 
The ``a'' in the lower caption is the intervention times as shown in Fig. \ref{fig:boid}. For example, ``$a = 95$'' means the case of intervention at the 95th frame (9.5 s).}
\textcolor{black}{The video is given in the above GitHub page.}
}
%
\label{fig:nba}
\end{figure}

%\vspace{-5pt}
\section{Conclusions}
\label{sec:conclusion}
In this paper, we proposed an interpretable counterfactual recurrent network in multiagent systems to estimate the effect of the intervention.
Using synthetic CARLA and Boid datasets, we showed that our model achieved lower errors in estimating counterfactual covariates and the most effective treatment timing than the baselines.
Furthermore, using real basketball data, our model performed realistic counterfactual prediction.
We consider a general ITE framework for various domains, but the experimental results show that for each domain the effective modeling was different. 
Possible future research directions are to realize better modeling of the future outcomes, 
% prediction of the potential outcomes 
and to apply our approach to other multiagent domains such as animals and pedestrians using domain knowledge.

\section*{Acknowledgments}
This work was supported by JSPS KAKENHI (Japan Society for the Promotion of Science, Grant Numbers 20H04075, 21H04892, and 21H05300), JST PRESTO (Japan Science and Technology Agency, Precursory Research for Embryonic Science and Technology, Grant Number JPMJPR20CA), and JST CREST (Core Research for Evolutional Science and Technology, Grant Number JPMJCR1913).
 

% An example of a floating figure using the graphicx package.
% Note that \label must occur AFTER (or within) \caption.
% For figures, \caption should occur after the \includegraphics.
% Note that IEEEtran v1.7 and later has special internal code that
% is designed to preserve the operation of \label within \caption
% even when the captionsoff option is in effect. However, because
% of issues like this, it may be the safest practice to put all your
% \label just after \caption rather than within \caption{}.
%
% Reminder: the "draftcls" or "draftclsnofoot", not "draft", class
% option should be used if it is desired that the figures are to be
% displayed while in draft mode.
%
%\begin{figure}[!t]
%\centering
%\includegraphics[width=2.5in]{myfigure}
% where an .eps filename suffix will be assumed under latex, 
% and a .pdf suffix will be assumed for pdflatex; or what has been declared
% via \DeclareGraphicsExtensions.
%\caption{Simulation results for the network.}
%\label{fig_sim}
%\end{figure}

% Note that the IEEE typically puts floats only at the top, even when this
% results in a large percentage of a column being occupied by floats.


% An example of a double column floating figure using two subfigures.
% (The subfig.sty package must be loaded for this to work.)
% The subfigure \label commands are set within each subfloat command,
% and the \label for the overall figure must come after \caption.
% \hfil is used as a separator to get equal spacing.
% Watch out that the combined width of all the subfigures on a 
% line do not exceed the text width or a line break will occur.
%
%\begin{figure*}[!t]
%\centering
%\subfloat[Case I]{\includegraphics[width=2.5in]{box}%
%\label{fig_first_case}}
%\hfil
%\subfloat[Case II]{\includegraphics[width=2.5in]{box}%
%\label{fig_second_case}}
%\caption{Simulation results for the network.}
%\label{fig_sim}
%\end{figure*}
%
% Note that often IEEE papers with subfigures do not employ subfigure
% captions (using the optional argument to \subfloat[]), but instead will
% reference/describe all of them (a), (b), etc., within the main caption.
% Be aware that for subfig.sty to generate the (a), (b), etc., subfigure
% labels, the optional argument to \subfloat must be present. If a
% subcaption is not desired, just leave its contents blank,
% e.g., \subfloat[].


% An example of a floating table. Note that, for IEEE style tables, the
% \caption command should come BEFORE the table and, given that table
% captions serve much like titles, are usually capitalized except for words
% such as a, an, and, as, at, but, by, for, in, nor, of, on, or, the, to
% and up, which are usually not capitalized unless they are the first or
% last word of the caption. Table text will default to \footnotesize as
% the IEEE normally uses this smaller font for tables.
% The \label must come after \caption as always.
%
%\begin{table}[!t]
%% increase table row spacing, adjust to taste
%\renewcommand{\arraystretch}{1.3}
% if using array.sty, it might be a good idea to tweak the value of
% \extrarowheight as needed to properly center the text within the cells
%\caption{An Example of a Table}
%\label{table_example}
%\centering
%% Some packages, such as MDW tools, offer better commands for making tables
%% than the plain LaTeX2e tabular which is used here.
%\begin{tabular}{|c||c|}
%\hline
%One & Two\\
%\hline
%Three & Four\\
%\hline
%\end{tabular}
%\end{table}


% Note that the IEEE does not put floats in the very first column
% - or typically anywhere on the first page for that matter. Also,
% in-text middle ("here") positioning is typically not used, but it
% is allowed and encouraged for Computer Society conferences (but
% not Computer Society journals). Most IEEE journals/conferences use
% top floats exclusively. 
% Note that, LaTeX2e, unlike IEEE journals/conferences, places
% footnotes above bottom floats. This can be corrected via the
% \fnbelowfloat command of the stfloats package.
\renewcommand{\thesection}{\Alph{section}}
%\appendix
\section*{Appendix}
%%%%%%%%%%%%%%%%%%%%%%%%%%%%%%%%%%%%%
\section{A proof of Theorem \ref{thm:identification}}
\label{app:proof}
\begin{proof}
Under the aforementioned assumptions in the main text, we can prove the identification of ITE:
\eq{\tau_t &= \mathbb{E}_y[ y_{1,t+1}-y_{0,t+1}|x_t,\mathcal{H}_t]\label{eq:proof1}\\
&=\mathbb{E}_z[\mathbb{E}_y[y_{1,t+1}-y_{0,t+1}|x_t,z_t,\mathcal{H}_t]|x_t,\mathcal{H}_t]\label{eq:proof2}\\
&=\mathbb{E}_z[\mathbb{E}_y[y_{1,t+1}-y_{0,t+1}|z_t]|x_t,\mathcal{H}_t]\label{eq:proof3}\\
&=\mathbb{E}_z[\mathbb{E}_y[y_{1,t+1}|z_t,a_t=1] - \mathbb{E}_y[y_{0,t+1}|z_t,a_t=0]|x_t,\mathcal{H}_t]\label{eq:proof4}\\
&=\mathbb{E}_z[\mathbb{E}_y[y_{F,t+1}|z_t,a_t=1] - \mathbb{E}_y[y_{F,t+1}|z_t,a_t=0]|x_t,\mathcal{H}_t],\label{eq:proof5}
%\label{eq:proof} 
}
where $\tau_t = \tau(x_t,\mathcal{H}_t)$, $y_{F,t+1}$ is a factual outcome, and we drop the instance index $(i)$ for simplification. 
Eq. (\ref{eq:proof1}) is the definition of ITE in our setting, Eq. (\ref{eq:proof2}) is a straightforward expectation over $p(z_t|x_t,\mathcal{H}_t)$ , and Eq. (\ref{eq:proof3}) be inferred from the structure of the causal graph shown in Fig. \ref{fig:causalgraph}.
%The SUTVA assumption is implicitly used in the causal graph, 
Eq. (\ref{eq:proof4}) is based on the assumption that $z_t$ contains all the hidden confounders, as well as the positivity assumption, and Eq. (\ref{eq:proof5}) can be inferred from the consistency assumption. 
Thus, if our framework can correctly model $p(z_t|x_t,\mathcal{H}_t)$ and $p(y_t|z_t,a_t)$, then the ITEs can be identified under the causal graph in Fig. \ref{fig:causalgraph}.
\end{proof}

%%%%%%%%%%%%%%%%%%%%%%%%%%%%%%%%%%%%%%
\section{Common training setup}
\label{app:common}
\subsection{Model training and computation}
\label{app:ourtraining}
The codes and data we used are provided at \url{https://github.com/keisuke198619/TGV-CRN}.
% via the supplementary materials. 
This experiment was performed on an Intel(R) Xeon(R) CPU E5-2699 v4 ($2.20$ GHz $\times$ 16) with GeForce TITAN X pascal GPU.
For the training of the proposed and baseline models, we used the Adam optimizer \cite{Kingma15} with an initial learning rate of $0.0001$ and $20$ training epochs.
We set the batchsize to 256.
For the hyper-parameters in the loss function, we set $\alpha = 0.1, \lambda = 0.1$ for all datasets and  $\gamma = 0.1$ in the Boid and CARLA dataset and  $\gamma = 1$ in the NBA experiment because the latter dataset was more difficult to predict than the Boid and CARLA experiments.
% The hidden layer is a two-layer MLPs of size ****. 

%In addition, to compare the methods in terms of their amount of computation, we measured training and inference time across two datasets.
% The results in Table \ref{tab:computation} show that the computation time of our method was...

\if0
\begin{table*}[ht!]
\centering
\scalebox{1}{
\begin{tabular}{l|cc}%|
\Xhline{3\arrayrulewidth} %\hline
& \me{1}{Kuramoto  model } & \me{1}{Boid model}  \\ 
& \me{1}{($p = 5, T = 200, K = 5$)}&\me{1}{($p = 5, T = 200, K = 3$)}\\
\hline
eSRU \cite{Khanna19} & 162 $\pm$ 5 & 143 $\pm$ 9 \\
GVAR \cite{Marcinkevics20} & 27 $\pm$ 3 & 19 $\pm$ 1
\\
\hline
ABM (full) & 116 $\pm$ 4 & 129 $\pm$ 4
\\
\Xhline{3\arrayrulewidth} % \hline
\end{tabular}
}
\caption{\label{tab:computation} The averaged computation time [s] among 10 sequences in two datasets.}
\end{table*}
\fi


\subsection{Baseline models implementation}
\label{app:baselines}
We compared the performances of our methods to infer ITE with those in the following baselines: a simple RNN using GRU \cite{Cho14},
deep sequential weighting (DSW) \cite{liu2020estimating}, dynamic networked observational data deconfounder (but for clarity, we change the name into GCRN: graph counterfactual recurrent network) \cite{ma2021deconfounding}.

\noindent {\bf{RNN}.} This approach is based on GRU \cite{Cho14}. This model predicts the input (covariates) at the next time stamp, the potential outcome, and the probability of receiving treatment.
We also model the hidden confounder as the hidden state of GRU, but do not learn the representation to reduce the confounding bias.  

\noindent {\bf{DSW}} \cite{liu2020estimating}. Compared with the original model, we modified it to our setting (e.g., multiagent, time-varying outcome, and long-term prediction), and to fairly compare with our model, we removed the attention module. 

\noindent {\bf{GCRN}} \cite{ma2021deconfounding}. Similarly, compared with the original model, we modified it to our setting (e.g., multiagent and long-term prediction) based on DSW, and to fairly compare with our model, we removed the attention module. 


%%%%%%%%%%%%%%%%%%%%%%%%%%%%%%%%%%%%%%
\section{Sensitivity analysis in hyperparameters}
\label{app:sensitivity}

We performed the sensitivity analysis in hyperparameters using the CARLA dataset. 
The hyperparameters are presented in Eq. (\ref{eq:loss}).
Results in Table \ref{tab:sensitivity} shows the existence of the trade-off between the prediction performances of the outcome and covariates ($\gamma$).  

\begin{table}[ht!]
\centering
\scalebox{0.75}{
\begin{tabular}{l|cccc}%|
\Xhline{3\arrayrulewidth} %\hline
% & \me{4}{carla simulation dataset}  \\ 
& \me{1}{$L_{Outcome}$ } & \me{1}{$\sqrt{\epsilon^t_{PEHE}}$} &  \me{1}{$\epsilon^t_{ATE}$} & \me{1}{$L_{Covariates}$} \\
\hline
$\alpha, \gamma, \lambda = 0.1$ (default) & 0.041 $\pm$ 0.003 &  0.020 $\pm$ 0.002 & 0.008 $\pm$ 0.001 & 0.098 $\pm$ 0.0004 \\
\hline
$\alpha = 1.0, \gamma, \lambda = 0.1$ & 0.040 $\pm$ 0.003 &  0.020 $\pm$ 0.002 & 0.009 $\pm$ 0.001 & 0.102 $\pm$ 0.0004 
\\$\alpha = 0.01, \gamma, \lambda = 0.1$ & 0.041 $\pm$ 0.003 &  0.019 $\pm$ 0.002 & 0.007 $\pm$ 0.001 & 0.097 $\pm$ 0.0004 
\\$\gamma = 1.0, \alpha, \lambda = 0.1$ & 0.041 $\pm$ 0.003 &  0.019 $\pm$ 0.001 & 0.005 $\pm$ 0.001 & 0.086 $\pm$ 0.0004\\
$\gamma = 0.01, \alpha, \lambda = 0.1$ &  0.031 $\pm$ 0.002 &  0.020 $\pm$ 0.002 & 0.009 $\pm$ 0.001 & 0.146 $\pm$ 0.0005  
\\$\lambda = 1.0, \alpha, \gamma = 0.1$ & 0.062 $\pm$ 0.005 &  0.019 $\pm$ 0.002 & 0.007 $\pm$ 0.001 & 0.121 $\pm$ 0.0005  
\\$\lambda = 0.01, \alpha, \gamma = 0.1$ &  0.041 $\pm$ 0.003 &  0.020 $\pm$ 0.002 & 0.008 $\pm$ 0.001 & 0.098 $\pm$ 0.0004
\\
\Xhline{3\arrayrulewidth} % \hline
\end{tabular}
}
\caption{\label{tab:sensitivity} Sensitivity analysis on the carla dataset.}
\end{table}

%%%%%%%%%%%%%%%%%%%%%%%%%%%%%%%%%%%%%%
\vspace{-10pt}
\section{Boid dataset}
\label{app:boid}
%\subsection{Simulation model}
% The rule-based models have been used to generate generic simulated flocking agents called boids \cite{reynolds1987flocks}.
The schooling model we used in this study was a unit-vector-based (rule-based) model \cite{Couzin02}, which accounts for the relative positions and direction vectors of neighboring fish agents, such that each fish tends to align its own direction vector with those of its neighbors. 
% In this paper, we intervene the agents' recognition to generate torus (circle) behaviors from the swarm (random) behaviors.
In this model, $20$ agents (length: 0.5 m) are described by a two-dimensional vector with a constant velocity (1 m/s) in a boundary square (30 $\times$ 30 m) as follows: 
${r}^k=\left({x_i}~{y_i}\right)^T$ and ${v}^k_t= \|v^k\|_2d_k$, where $x_i$ and $y_i$ are two-dimensional Cartesian coordinates, ${v}^k$ is a velocity vector, $\|\cdot\|_2$ is the Euclidean norm, and $d_k$ is an unit directional vector for agent $i$.

At each timestep, a member will change direction according to the positions of all other members. The space around an individual is divided into three zones where each modifying the unit vector of the velocity (for the zones, see the main text).
%The first region, called the repulsion zone with radius $r_r = 0.5$ m, corresponds to the ``personal'' space of the particle. Individuals within each other’s repulsion zones will try to avoid each other by swimming in opposite directions. 
%The second region is called the orientation zone, in which members try to move in the same direction  (radius $r_o$). 
%We set $r_o = 4$ to generate swarming behaviors before the intervention. 
%To generate torus behaviors, we change $r_o = 5$, which is the intervention in this experiment.
%The third is the attractive zone (radius $r_a = 7.5$ m), in which agents move towards each other and tend to cluster, while any agents beyond that radius have no influence. 
Let $\lambda_r$, $\lambda_o$, and $\lambda_a$ be the numbers in the zones of repulsion, orientation and attraction respectively. For $\lambda_r \neq 0$, the unit vector of an individual at the next timestep is given by:
\begin{equation} 
\label{eq:boid1}
d_k(t+1 , \lambda_r \neq 0 )=-\left(\frac{1}{\lambda_r-1}\sum_{j\neq k}^{\lambda_r}\frac{r^{kj}_t}{\|r^{kj}_t\|_2}\right),
\end{equation} 
where $r^{kj}={r}_j-{r}_i$.
The velocity vector points away from neighbors within this zone to prevent collisions. This zone is given the highest priority; if and only if $\lambda_r = 0$, the remaining zones are considered. 
The unit vector in this case is given by:
\begin{equation} 
\label{eq:boid2}
d_k(t+1 , {\lambda}_r=0)=\frac{1}{2}\left(\frac{1}{\lambda_o}\sum_{j=1}^{\lambda_o} d_j(t)+\frac{1}{\lambda_a-1}\sum_{j\neq k}^{{\lambda}_a}\frac{r^{kj}_t}{\|r^{kj}_t\|_2}\right).
\end{equation} 
The first term corresponds to the orientation zone while the second term corresponds to the attraction zone. The above equation contains a factor of $1/2$ which normalizes the unit vector in the case where both zones have non-zero neighbors. If no agents are found near any zone, the individual maintains a constant velocity at each timestep.

In addition to the above, we constrain the angle by which a member can change its unit vector at each timestep to a maximum of $\beta = 30$ deg. This condition was imposed to facilitate rigid body dynamics. Since we assumed point-like members, all information about the physical dimensions of the actual fish is lost, which leaves the unit vector free to rotate at any angle. In reality, however, the conservation of angular momentum will limit the ability of the fish to turn angle $\theta$ as follows:
\begin{equation} 
\label{eq:boid3}
  d_k\left(t+1 \right)\cdot d_k(t) = 
  \begin{cases}
   \cos(\beta ) & \text{if $\theta >\beta$} \\
   \cos\left(\theta \right) & \text{otherwise}.
  \end{cases}
\end{equation} 
If the above condition is not satisfied, the angle of the desired direction at the next timestep is rescaled to $\theta = \beta$. In this way, any un-physical behavior such as having a 180$^\circ$ rotation of the velocity vector in a single timestep is prevented.

In the simulation, the ground truth of $\tau_T$ was $0.091 \pm 0.026$, which indicates the intervention increased angular velocities (see also the main text). 


% \section{Results of the boid model}
% \label{app:res_boid}

%%%%%%%%%%%%%%%%%%%%%%%%%%%%%%%%%%%%%%
\section{CARLA dataset}
\label{app:carla}
%\subsection{Simulation model}
We used the CARLA simulator \cite{dosovitskiy2017carla} (ver. 0.9.8).
The CARLA environment contains dynamic obstacles (e.g., pedestrians and cars) that interact with the ego car. For generating data, we performed simulation through various towns, starting points, and obstacle positions, which were subsampled at 2 Hz.
The obstacles' positions and starting points of the ego car were randomly selected from the possible locations on the map for each run.
Approximately 20-120 obstacles were placed on each map.
Using the data collected in CARLA, the same driving conditions were reproduced in ROS (robot operating system) \cite{quigley2009ros}, and Autoware was used to make the vehicle drive the same route autonomously.
Here we consider two types of autonomous vehicles: partial observation with Autoware \cite{kato2015open} (ver.1.12.0) and full observation types with the pre-installed CARLA simulator. 
The partial observation model often stops when dangerous situations for safety. 
If an obstacle enters the deceleration or stopping range, the vehicle decelerates or stops in front of the obstacle.
% The yellow area in Fig.\ref{fig:intervention} is the deceleration range, and the green area is the stopping range. 
% If an obstacle enters the area, the vehicle decelerates or stops in front of the obstacle.
The deceleration range was set to be wider than usual for safety.
The full observation model intervenes to accelerate the stopped ego car after the safety confirmation based on the speed information at the data collection.

We set $T=60, T_b=40$, and $T_i=T_b+1,\ldots,T_b+10$.
% \tau = 1$.
% For the counterfactual data, since the dangerous situations where intervention is required are limited, in this experiment we only used with or without intervention at the same timing (not various timings). 
We evaluated and predicted the safe driving distance of the ego car from the starting point while driving without collision with any obstacle (for details, see the main text). 
% That is, the treatment effect of the full observation model is defined as the safe driving distance.
In the simulation, the ground truth of $\tau_T$ was $0.013 \pm 0.001$, which indicates the intervention increased safe driving distance. 
%%%%%%%%%%%%%%%%%%%%%%%%%%%%%%%%%%%%%%%
\section{NBA dataset}
\label{app:nba}
Here, we describe the details of the computation using the NBA dataset.
Dataset description is given in the main text.
We describe the computation of the attack effectiveness used in this study.
Then, we predict the effective attacks at the next time stamp using the pre-training samples and a logistic regression.  %  (i.e., $\tau=1$)
%
% It could be argued that the tactics and strategy of a coach and team are perhaps most influential up until the point at which there is a good scoring opportunity to make a shot, and it is then up to the skills/form of the individual player to actually convert this opportunity into points.
%
% A good scoring opportunity in basketball is considered to be a shot being made where there has been a high expected probability of scoring in the past.
%

%
From the aforementioned reasons in the main text, we compute an interpretable and simple indicator from available statistics (i.e., based on the frequency) to evaluate whether a player attempts a better shot, rather than based on the shot label or learning-based score prediction.
From available statistics, we focused on two basic factors for effective attacks at an individual player level: the shot zone on the court and the distance between a shooter and the nearest defender. 
These two factors have been considered to be important for basketball shot prediction \cite{Fujii16,Fujii17,Fujii18}.
In the NBA advanced stats (\url{https://www.nba.com/stats/players/shots-closest-defender/}), we can access probabilities of successful shots in each zone and distance for each player.
The shot zones are separated into four areas: the restricted area, in-the-paint, mid-range, and the 3-point area. 
The restricted area is defined as the area within a radius of 2.44 m from the ring (distance between the side of the rectangle and the ring) from the ring.
In-the-paint is defined as the area within a radius of 5.46 m (distance between the ring and the farthest vertex of the rectangle) from the ring.
The 3-point area is defined as the outside of the 3-point line. 
Mid-range is the remaining area. 
The shooter's distance from the nearest defender is categorized into four ranges: $0-2$ feet, $2-4$ feet, $4-6$ feet, and $6+$ feet.

We define the attack effectiveness using the following criteria:
\vspace{-0pt}
\begin{quote}
\begin{itemize}
\item The shooter's position in the restricted area is effective at any distance because the defender often exists near the shooter.
\item The shooter's position in the paint and mid-range is effective at 6 feet or further (this range is regarded as ``open'' in the NBA advanced stats).
\item The shooter's position in the 3-point area is effective when a player with a probability of 0.35 and hits with 6 feet or more (because some players do not shoot tactically).
\end{itemize}
\end{quote}
\vspace{-0pt}
Based on the statistics before the game (e.g., for the pre-training data, we used the 2014/2015 season statistics) and the tracking data, 
we computed the probabilities of successful shots for each zone and the distances for each player. 
We computed the probability of the player who attempted shots less than 10 times as the probability of the same position player (i.e., guard, forward, center, guard/forward, and forward/center based on the registration information from the NBA 2014/2015 season).
It should be noted that there are some strategies of a good shot in basketball that differ depending on the court location and context, for example, 2 pointers and 3 pointers.
Note that, unfortunately, we can access those for only two areas (2-point and 3-point areas) with four distance categories, thus we computed the shot success probability at the restricted area, in-the-paint, and mid-range using that at the 2-point area.
% For adjusting the 3-point shot probability far from the 3-point line, we adjusted the probability such that the probability is linearly reduced by 0.2 at 12.73m (distance between the ring and the half-court line).
%
Based on the above definitions, for pre-training data, there were 13,681 shot successes, 15,443 shot failures, 5,572 turnovers,
16,976 effective attacks, and 17,720 ineffective attacks.
For training data, there were 4,677 shot successes, 5,092 shot failures, 1,691 turnovers, 5,602 effective attacks, and 5,858 ineffective attacks.
For test data, there were 517 shot successes, 577 shot failures, 211 turnovers, 664 effective attacks, and 641 ineffective attacks.
The probabilities of scoring, given the attack was effective and ineffective, were 0.463/0.415/0.417 and 0.328/0.402/0.374 for pre-training, training, and test data, respectively.
We confirmed that the effective attack had a higher probability of a successful shot than the ineffective attack. 

After the computation of the attack effectiveness, we predict the effective attacks at the next time stamp using the pre-training samples and a logistic regression.  
In the soccer domain, there has been some work to evaluate players and teams (e.g., \cite{Decroos19,Toda2021}) based on the prediction model (i.e., classifier) on the assumption that a good play is a play that will bring good outcomes (e.g., score and ball recovery) in the future. 
These approaches can transform a discrete value (i.e., an outcome) into a continuous value (e.g., a probability).
According to these papers, we also predict the effective attack at the next time stamp.
To verify the prediction accuracy, we compared the accuracy of the logistic regression, LightGBM \cite{ke2017lightgbm}, which is a popular classifier using a highly efficient gradient boosting decision tree, and prediction with all effective attacks.
For the training data, the accuracies of logistic regression, LightGBM, and the prediction with all the same labels were 0.838, 0.838, and 0.731, respectively.
For the test data, these were 0.679, 0.676, and 0.568, respectively. 
We confirmed that the logistic regression had competitive performance compared with LightGBM, and higher accuracy than the prediction with all the same labels, even maintaining higher interpretability. 


\section{\textcolor{black}{Additional analysis in basketball dataset}}
\label{app:analysis_basket}
\textcolor{black}{As additional analysis, we indicate endpoint errors and distributions of long-term covariate prediction for basketball data.
Results in Table \ref{tab:nba_end} that the tendency of the endpoint prediction error for all models is similar to the mean prediction error, in which our approach (TV-CRN and TGV-CRN) show the best performance. In addition, the distribution of the player velocity in Fig. \ref{fig:histogram_nba} shows that our approach without theory-based computation had a wider distribution like ground truth than other models, and those with theory-based computation had less zero-velocity bins like ground truth than other models. }
\begin{table}[ht!]
\centering
\scalebox{1}{
\begin{tabular}{l|c}%|
\Xhline{3\arrayrulewidth} %\hline
% & \me{4}{carla simulation dataset}  \\ 
 & \me{1}{$L_{Covariates}$ at endpoint} \\
\hline
RNN & 1.637 $\pm$ 0.0096
\\DSW \cite{liu2020estimating}& ---
\\GCRN \cite{ma2021deconfounding} & ---
\\GCRN$+$X & 2.266 $\pm$ 0.0138\\
\hline
GV-CRN & 1.937 $\pm$ 0.0116
\\TG-CRN & \textbf{1.482} $\pm$ \textbf{0.0090}
\\TV-CRN & 1.566 $\pm$ 0.0094
\\ TGV-CRN (full) & \textbf{ 1.510} $\pm$ \textbf{0.0093}
\\
\Xhline{3\arrayrulewidth} % \hline
\end{tabular}
}
\vspace{3pt}
\caption{\label{tab:nba_end} \textcolor{black}{Endpoint prediction error on NBA dataset.}}
\end{table}

\begin{figure}[h]
\centering
\includegraphics[width=0.5\textwidth]{fig/histogram_nba.png}
\caption[]{\textcolor{black}{Histograms of player velocity in NBA dataset.}
}
%
\label{fig:histogram_nba}
\end{figure}


%\section{GRL illustration}
%\label{app:GRL}
\section{\textcolor{black}{Convergence in the learning of our models}}
\label{app:converge}
\textcolor{black}{We illustrate the change in the validation losses of our models over the course of training epochs as shown in Figs. \ref{fig:carla_learning}, \ref{fig:boid_learning}, and \ref{fig:nba_learning}.
Compared with CARLA dataset results, which show better outcome and covariate prediction results than other datasets, those in Boid and NBA datasets sometimes show unstable learning curves, but most of the models finally show convergence in all datasets. 
To perform fair comparisons among the models, we trained all models with 20 epochs. }

\begin{figure}[h]
\centering
\includegraphics[width=0.5\textwidth]{fig/carla_learning.png}
\caption[]{\textcolor{black}{The learning curve of our model training in CARLA dataset. Left and right subfigures indicate losses in outcome and covariate prediction, respectively. }
}
%
\label{fig:carla_learning}
\end{figure}

\begin{figure}[h]
\centering
\includegraphics[width=0.5\textwidth]{fig/boid_learning.png}
\caption[]{\textcolor{black}{The learning curve of our model training in Boid dataset.}
}
%
\label{fig:boid_learning}
\end{figure}

\begin{figure}[h]
\centering
\includegraphics[width=0.5\textwidth]{fig/nba_learning.png}
\caption[]{\textcolor{black}{The learning curve of our model training in NBA dataset.}
}
%
\label{fig:nba_learning}
\end{figure}
% if have a single appendix:
%\appendix[Proof of the Zonklar Equations]
% or
%\appendix  % for no appendix heading
% do not use \section anymore after \appendix, only \section*
% is possibly needed

% use appendices with more than one appendix
% then use \section to start each appendix
% you must declare a \section before using any
% \subsection or using \label (\appendices by itself
% starts a section numbered zero.)
%
\bibliographystyle{IEEEtran}
\bibliography{main}

\clearpage
% \renewcommand{\thesection}{\Alph{section}}
%\appendix
\section*{Appendix}
%%%%%%%%%%%%%%%%%%%%%%%%%%%%%%%%%%%%%
\section{A proof of Theorem \ref{thm:identification}}
\label{app:proof}
\begin{proof}
Under the aforementioned assumptions in the main text, we can prove the identification of ITE:
\eq{\tau_t &= \mathbb{E}_y[ y_{1,t+1}-y_{0,t+1}|x_t,\mathcal{H}_t]\label{eq:proof1}\\
&=\mathbb{E}_z[\mathbb{E}_y[y_{1,t+1}-y_{0,t+1}|x_t,z_t,\mathcal{H}_t]|x_t,\mathcal{H}_t]\label{eq:proof2}\\
&=\mathbb{E}_z[\mathbb{E}_y[y_{1,t+1}-y_{0,t+1}|z_t]|x_t,\mathcal{H}_t]\label{eq:proof3}\\
&=\mathbb{E}_z[\mathbb{E}_y[y_{1,t+1}|z_t,a_t=1] - \mathbb{E}_y[y_{0,t+1}|z_t,a_t=0]|x_t,\mathcal{H}_t]\label{eq:proof4}\\
&=\mathbb{E}_z[\mathbb{E}_y[y_{F,t+1}|z_t,a_t=1] - \mathbb{E}_y[y_{F,t+1}|z_t,a_t=0]|x_t,\mathcal{H}_t],\label{eq:proof5}
%\label{eq:proof} 
}
where $\tau_t = \tau(x_t,\mathcal{H}_t)$, $y_{F,t+1}$ is a factual outcome, and we drop the instance index $(i)$ for simplification. 
Eq. (\ref{eq:proof1}) is the definition of ITE in our setting, Eq. (\ref{eq:proof2}) is a straightforward expectation over $p(z_t|x_t,\mathcal{H}_t)$ , and Eq. (\ref{eq:proof3}) be inferred from the structure of the causal graph shown in Fig. \ref{fig:causalgraph}.
%The SUTVA assumption is implicitly used in the causal graph, 
Eq. (\ref{eq:proof4}) is based on the assumption that $z_t$ contains all the hidden confounders, as well as the positivity assumption, and Eq. (\ref{eq:proof5}) can be inferred from the consistency assumption. 
Thus, if our framework can correctly model $p(z_t|x_t,\mathcal{H}_t)$ and $p(y_t|z_t,a_t)$, then the ITEs can be identified under the causal graph in Fig. \ref{fig:causalgraph}.
\end{proof}

%%%%%%%%%%%%%%%%%%%%%%%%%%%%%%%%%%%%%%
\section{Common training setup}
\label{app:common}
\subsection{Model training and computation}
\label{app:ourtraining}
The codes and data we used are provided at \url{https://github.com/keisuke198619/TGV-CRN}.
% via the supplementary materials. 
This experiment was performed on an Intel(R) Xeon(R) CPU E5-2699 v4 ($2.20$ GHz $\times$ 16) with GeForce TITAN X pascal GPU.
For the training of the proposed and baseline models, we used the Adam optimizer \cite{Kingma15} with an initial learning rate of $0.0001$ and $20$ training epochs.
We set the batchsize to 256.
For the hyper-parameters in the loss function, we set $\alpha = 0.1, \lambda = 0.1$ for all datasets and  $\gamma = 0.1$ in the Boid and CARLA dataset and  $\gamma = 1$ in the NBA experiment because the latter dataset was more difficult to predict than the Boid and CARLA experiments.
% The hidden layer is a two-layer MLPs of size ****. 

%In addition, to compare the methods in terms of their amount of computation, we measured training and inference time across two datasets.
% The results in Table \ref{tab:computation} show that the computation time of our method was...

\if0
\begin{table*}[ht!]
\centering
\scalebox{1}{
\begin{tabular}{l|cc}%|
\Xhline{3\arrayrulewidth} %\hline
& \me{1}{Kuramoto  model } & \me{1}{Boid model}  \\ 
& \me{1}{($p = 5, T = 200, K = 5$)}&\me{1}{($p = 5, T = 200, K = 3$)}\\
\hline
eSRU \cite{Khanna19} & 162 $\pm$ 5 & 143 $\pm$ 9 \\
GVAR \cite{Marcinkevics20} & 27 $\pm$ 3 & 19 $\pm$ 1
\\
\hline
ABM (full) & 116 $\pm$ 4 & 129 $\pm$ 4
\\
\Xhline{3\arrayrulewidth} % \hline
\end{tabular}
}
\caption{\label{tab:computation} The averaged computation time [s] among 10 sequences in two datasets.}
\end{table*}
\fi


\subsection{Baseline models implementation}
\label{app:baselines}
We compared the performances of our methods to infer ITE with those in the following baselines: a simple RNN using GRU \cite{Cho14},
deep sequential weighting (DSW) \cite{liu2020estimating}, dynamic networked observational data deconfounder (but for clarity, we change the name into GCRN: graph counterfactual recurrent network) \cite{ma2021deconfounding}.

\noindent {\bf{RNN}.} This approach is based on GRU \cite{Cho14}. This model predicts the input (covariates) at the next time stamp, the potential outcome, and the probability of receiving treatment.
We also model the hidden confounder as the hidden state of GRU, but do not learn the representation to reduce the confounding bias.  

\noindent {\bf{DSW}} \cite{liu2020estimating}. Compared with the original model, we modified it to our setting (e.g., multi-agent, time-varying outcome, and long-term prediction), and to fairly compare with our model, we removed the attention module. 

\noindent {\bf{GCRN}} \cite{ma2021deconfounding}. Similarly, compared with the original model, we modified it to our setting (e.g., multi-agent and long-term prediction) based on DSW, and to fairly compare with our model, we removed the attention module. 


%%%%%%%%%%%%%%%%%%%%%%%%%%%%%%%%%%%%%%
\section{Sensitivity analysis in hyperparameters}
\label{app:sensitivity}

We performed the sensitivity analysis in hyperparameters using the CARLA dataset. 
The hyperparameters are presented in Eq. (\ref{eq:loss}).
Results in Table \ref{tab:sensitivity} shows the existence of the trade-off between the prediction performances of the outcome and covariates ($\gamma$).  

\begin{table}[ht!]
\centering
\scalebox{0.75}{
\begin{tabular}{l|cccc}%|
\Xhline{3\arrayrulewidth} %\hline
% & \me{4}{carla simulation dataset}  \\ 
& \me{1}{$L_{Outcome}$ } & \me{1}{$\sqrt{\epsilon^t_{PEHE}}$} &  \me{1}{$\epsilon^t_{ATE}$} & \me{1}{$L_{Covariates}$} \\
\hline
$\alpha, \gamma, \lambda = 0.1$ (default) & 0.041 $\pm$ 0.003 &  0.020 $\pm$ 0.002 & 0.008 $\pm$ 0.001 & 0.098 $\pm$ 0.0004 \\
\hline
$\alpha = 1.0, \gamma, \lambda = 0.1$ & 0.040 $\pm$ 0.003 &  0.020 $\pm$ 0.002 & 0.009 $\pm$ 0.001 & 0.102 $\pm$ 0.0004 
\\$\alpha = 0.01, \gamma, \lambda = 0.1$ & 0.041 $\pm$ 0.003 &  0.019 $\pm$ 0.002 & 0.007 $\pm$ 0.001 & 0.097 $\pm$ 0.0004 
\\$\gamma = 1.0, \alpha, \lambda = 0.1$ & 0.041 $\pm$ 0.003 &  0.019 $\pm$ 0.001 & 0.005 $\pm$ 0.001 & 0.086 $\pm$ 0.0004\\
$\gamma = 0.01, \alpha, \lambda = 0.1$ &  0.031 $\pm$ 0.002 &  0.020 $\pm$ 0.002 & 0.009 $\pm$ 0.001 & 0.146 $\pm$ 0.0005  
\\$\lambda = 1.0, \alpha, \gamma = 0.1$ & 0.062 $\pm$ 0.005 &  0.019 $\pm$ 0.002 & 0.007 $\pm$ 0.001 & 0.121 $\pm$ 0.0005  
\\$\lambda = 0.01, \alpha, \gamma = 0.1$ &  0.041 $\pm$ 0.003 &  0.020 $\pm$ 0.002 & 0.008 $\pm$ 0.001 & 0.098 $\pm$ 0.0004
\\
\Xhline{3\arrayrulewidth} % \hline
\end{tabular}
}
\caption{\label{tab:sensitivity} Sensitivity analysis on the carla dataset.}
\end{table}

%%%%%%%%%%%%%%%%%%%%%%%%%%%%%%%%%%%%%%
\vspace{-10pt}
\section{Boid dataset}
\label{app:boid}
%\subsection{Simulation model}
% The rule-based models have been used to generate generic simulated flocking agents called boids \cite{reynolds1987flocks}.
The schooling model we used in this study was a unit-vector-based (rule-based) model \cite{Couzin02}, which accounts for the relative positions and direction vectors of neighboring fish agents, such that each fish tends to align its own direction vector with those of its neighbors. 
% In this paper, we intervene the agents' recognition to generate torus (circle) behaviors from the swarm (random) behaviors.
In this model, $20$ agents (length: 0.5 m) are described by a two-dimensional vector with a constant velocity (1 m/s) in a boundary square (30 $\times$ 30 m) as follows: 
${r}^k=\left({x_i}~{y_i}\right)^T$ and ${v}^k_t= \|v^k\|_2d_k$, where $x_i$ and $y_i$ are two-dimensional Cartesian coordinates, ${v}^k$ is a velocity vector, $\|\cdot\|_2$ is the Euclidean norm, and $d_k$ is an unit directional vector for agent $i$.

At each timestep, a member will change direction according to the positions of all other members. The space around an individual is divided into three zones where each modifying the unit vector of the velocity (for the zones, see the main text).
%The first region, called the repulsion zone with radius $r_r = 0.5$ m, corresponds to the ``personal'' space of the particle. Individuals within each other’s repulsion zones will try to avoid each other by swimming in opposite directions. 
%The second region is called the orientation zone, in which members try to move in the same direction  (radius $r_o$). 
%We set $r_o = 4$ to generate swarming behaviors before the intervention. 
%To generate torus behaviors, we change $r_o = 5$, which is the intervention in this experiment.
%The third is the attractive zone (radius $r_a = 7.5$ m), in which agents move towards each other and tend to cluster, while any agents beyond that radius have no influence. 
Let $\lambda_r$, $\lambda_o$, and $\lambda_a$ be the numbers in the zones of repulsion, orientation and attraction respectively. For $\lambda_r \neq 0$, the unit vector of an individual at the next timestep is given by:
\begin{equation} 
\label{eq:boid1}
d_k(t+1 , \lambda_r \neq 0 )=-\left(\frac{1}{\lambda_r-1}\sum_{j\neq k}^{\lambda_r}\frac{r^{kj}_t}{\|r^{kj}_t\|_2}\right),
\end{equation} 
where $r^{kj}={r}_j-{r}_i$.
The velocity vector points away from neighbors within this zone to prevent collisions. This zone is given the highest priority; if and only if $\lambda_r = 0$, the remaining zones are considered. 
The unit vector in this case is given by:
\begin{equation} 
\label{eq:boid2}
d_k(t+1 , {\lambda}_r=0)=\frac{1}{2}\left(\frac{1}{\lambda_o}\sum_{j=1}^{\lambda_o} d_j(t)+\frac{1}{\lambda_a-1}\sum_{j\neq k}^{{\lambda}_a}\frac{r^{kj}_t}{\|r^{kj}_t\|_2}\right).
\end{equation} 
The first term corresponds to the orientation zone while the second term corresponds to the attraction zone. The above equation contains a factor of $1/2$ which normalizes the unit vector in the case where both zones have non-zero neighbors. If no agents are found near any zone, the individual maintains a constant velocity at each timestep.

In addition to the above, we constrain the angle by which a member can change its unit vector at each timestep to a maximum of $\beta = 30$ deg. This condition was imposed to facilitate rigid body dynamics. Since we assumed point-like members, all information about the physical dimensions of the actual fish is lost, which leaves the unit vector free to rotate at any angle. In reality, however, the conservation of angular momentum will limit the ability of the fish to turn angle $\theta$ as follows:
\begin{equation} 
\label{eq:boid3}
  d_k\left(t+1 \right)\cdot d_k(t) = 
  \begin{cases}
   \cos(\beta ) & \text{if $\theta >\beta$} \\
   \cos\left(\theta \right) & \text{otherwise}.
  \end{cases}
\end{equation} 
If the above condition is not satisfied, the angle of the desired direction at the next timestep is rescaled to $\theta = \beta$. In this way, any un-physical behavior such as having a 180$^\circ$ rotation of the velocity vector in a single timestep is prevented.

In the simulation, the ground truth of $\tau_T$ was $0.091 \pm 0.026$, which indicates the intervention increased angular velocities (see also the main text). 


% \section{Results of the boid model}
% \label{app:res_boid}

%%%%%%%%%%%%%%%%%%%%%%%%%%%%%%%%%%%%%%
\section{CARLA dataset}
\label{app:carla}
%\subsection{Simulation model}
We used the CARLA simulator \cite{dosovitskiy2017carla} (ver. 0.9.8).
The CARLA environment contains dynamic obstacles (e.g., pedestrians and cars) that interact with the ego car. For generating data, we performed simulation through various towns, starting points, and obstacle positions, which were subsampled at 2 Hz.
The obstacles' positions and starting points of the ego car were randomly selected from the possible locations on the map for each run.
Approximately 20-120 obstacles were placed on each map.
Using the data collected in CARLA, the same driving conditions were reproduced in ROS (robot operating system) \cite{quigley2009ros}, and Autoware was used to make the vehicle drive the same route autonomously.
Here we consider two types of autonomous vehicles: partial observation with Autoware \cite{kato2015open} (ver.1.12.0) and full observation types with the pre-installed CARLA simulator. 
The partial observation model often stops when dangerous situations for safety. 
If an obstacle enters the deceleration or stopping range, the vehicle decelerates or stops in front of the obstacle.
% The yellow area in Fig.\ref{fig:intervention} is the deceleration range, and the green area is the stopping range. 
% If an obstacle enters the area, the vehicle decelerates or stops in front of the obstacle.
The deceleration range was set to be wider than usual for safety.
The full observation model intervenes to accelerate the stopped ego car after the safety confirmation based on the speed information at the data collection.

We set $T=60, T_b=40$, and $T_i=T_b+1,\ldots,T_b+10$.
% \tau = 1$.
% For the counterfactual data, since the dangerous situations where intervention is required are limited, in this experiment we only used with or without intervention at the same timing (not various timings). 
We evaluated and predicted the safe driving distance of the ego car from the starting point while driving without collision with any obstacle (for details, see the main text). 
% That is, the treatment effect of the full observation model is defined as the safe driving distance.
In the simulation, the ground truth of $\tau_T$ was $0.013 \pm 0.001$, which indicates the intervention increased safe driving distance. 
%%%%%%%%%%%%%%%%%%%%%%%%%%%%%%%%%%%%%%%
\section{NBA dataset}
\label{app:nba}
Here, we describe the details of the computation using the NBA dataset.
Dataset description is given in the main text.
We describe the computation of the attack effectiveness used in this study.
Then, we predict the effective attacks at the next time stamp using the pre-training samples and a logistic regression.  %  (i.e., $\tau=1$)
%
% It could be argued that the tactics and strategy of a coach and team are perhaps most influential up until the point at which there is a good scoring opportunity to make a shot, and it is then up to the skills/form of the individual player to actually convert this opportunity into points.
%
% A good scoring opportunity in basketball is considered to be a shot being made where there has been a high expected probability of scoring in the past.
%

%
From the aforementioned reasons in the main text, we compute an interpretable and simple indicator from available statistics (i.e., based on the frequency) to evaluate whether a player attempts a better shot, rather than based on the shot label or learning-based score prediction.
From available statistics, we focused on two basic factors for effective attacks at an individual player level: the shot zone on the court and the distance between a shooter and the nearest defender. 
These two factors have been considered to be important for basketball shot prediction \cite{Fujii16,Fujii17,Fujii18}.
In the NBA advanced stats (\url{https://www.nba.com/stats/players/shots-closest-defender/}), we can access probabilities of successful shots in each zone and distance for each player.
The shot zones are separated into four areas: the restricted area, in-the-paint, mid-range, and the 3-point area. 
The restricted area is defined as the area within a radius of 2.44 m from the ring (distance between the side of the rectangle and the ring) from the ring.
In-the-paint is defined as the area within a radius of 5.46 m (distance between the ring and the farthest vertex of the rectangle) from the ring.
The 3-point area is defined as the outside of the 3-point line. 
Mid-range is the remaining area. 
The shooter's distance from the nearest defender is categorized into four ranges: $0-2$ feet, $2-4$ feet, $4-6$ feet, and $6+$ feet.

We define the attack effectiveness using the following criteria:
\vspace{-0pt}
\begin{quote}
\begin{itemize}
\item The shooter's position in the restricted area is effective at any distance because the defender often exists near the shooter.
\item The shooter's position in the paint and mid-range is effective at 6 feet or further (this range is regarded as ``open'' in the NBA advanced stats).
\item The shooter's position in the 3-point area is effective when a player with a probability of 0.35 and hits with 6 feet or more (because some players do not shoot tactically).
\end{itemize}
\end{quote}
\vspace{-0pt}
Based on the statistics before the game (e.g., for the pre-training data, we used the 2014/2015 season statistics) and the tracking data, 
we computed the probabilities of successful shots for each zone and the distances for each player. 
We computed the probability of the player who attempted shots less than 10 times as the probability of the same position player (i.e., guard, forward, center, guard/forward, and forward/center based on the registration information from the NBA 2014/2015 season).
It should be noted that there are some strategies of a good shot in basketball that differ depending on the court location and context, for example, 2 pointers and 3 pointers.
Note that, unfortunately, we can access those for only two areas (2-point and 3-point areas) with four distance categories, thus we computed the shot success probability at the restricted area, in-the-paint, and mid-range using that at the 2-point area.
% For adjusting the 3-point shot probability far from the 3-point line, we adjusted the probability such that the probability is linearly reduced by 0.2 at 12.73m (distance between the ring and the half-court line).
%
Based on the above definitions, for pre-training data, there were 13,681 shot successes, 15,443 shot failures, 5,572 turnovers,
16,976 effective attacks, and 17,720 ineffective attacks.
For training data, there were 4,677 shot successes, 5,092 shot failures, 1,691 turnovers, 5,602 effective attacks, and 5,858 ineffective attacks.
For test data, there were 517 shot successes, 577 shot failures, 211 turnovers, 664 effective attacks, and 641 ineffective attacks.
The probabilities of scoring, given the attack was effective and ineffective, were 0.463/0.415/0.417 and 0.328/0.402/0.374 for pre-training, training, and test data, respectively.
We confirmed that the effective attack had a higher probability of a successful shot than the ineffective attack. 

After the computation of the attack effectiveness, we predict the effective attacks at the next time stamp using the pre-training samples and a logistic regression.  
In the soccer domain, there has been some work to evaluate players and teams (e.g., \cite{Decroos19,Toda2021}) based on the prediction model (i.e., classifier) on the assumption that a good play is a play that will bring good outcomes (e.g., score and ball recovery) in the future. 
These approaches can transform a discrete value (i.e., an outcome) into a continuous value (e.g., a probability).
According to these papers, we also predict the effective attack at the next time stamp.
To verify the prediction accuracy, we compared the accuracy of the logistic regression, LightGBM \cite{ke2017lightgbm}, which is a popular classifier using a highly efficient gradient boosting decision tree, and prediction with all effective attacks.
For the training data, the accuracies of logistic regression, LightGBM, and the prediction with all the same labels were 0.838, 0.838, and 0.731, respectively.
For the test data, these were 0.679, 0.676, and 0.568, respectively. 
We confirmed that the logistic regression had competitive performance compared with LightGBM, and higher accuracy than the prediction with all the same labels, even maintaining higher interpretability. 

%\section{GRL illustration}
%\label{app:GRL}
\section{\textcolor{black}{Convergence in the learning of our models}}
\label{app:converge}
\textcolor{black}{We illustrate the change in the validation losses of our models over the course of training epochs as shown in Figs. \ref{fig:carla_learning}, \ref{fig:boid_learning}, and \ref{fig:nba_learning}.
Compared with CARLA dataset results, which show better outcome and covariate prediction results than other datasets, those in Boid and NBA datasets sometimes show unstable learning curves, but most of the models finally show convergence in all datasets. 
To perform fair comparisons among the models, we trained all models with 20 epochs. }

\begin{figure}[h]
\centering
\includegraphics[width=0.5\textwidth]{fig/carla_learning.png}
\caption[]{\textcolor{black}{The learning curve of our model training in CARLA dataset. Left and right subfigures indicate losses in outcome and covariate prediction, respectively. }
}
%
\label{fig:carla_learning}
\end{figure}

\begin{figure}[h]
\centering
\includegraphics[width=0.5\textwidth]{fig/boid_learning.png}
\caption[]{\textcolor{black}{The learning curve of our model training in Boid dataset.}
}
%
\label{fig:boid_learning}
\end{figure}

\begin{figure}[h]
\centering
\includegraphics[width=0.5\textwidth]{fig/nba_learning.png}
\caption[]{\textcolor{black}{The learning curve of our model training in NBA dataset.}
}
%
\label{fig:nba_learning}
\end{figure}

\section{\textcolor{black}{Additional analysis in basketball dataset}}
\label{app:analysis_basket}
\textcolor{black}{As additional analysis, we indicate endpoint errors and distributions of long-term covariate prediction for basketball data.
Results in Table \ref{tab:nba_end} that the tendency of the endpoint prediction error for all models is similar to the mean prediction error, in which our approach (TV-CRN and TGV-CRN) show the best performance. In addition, the distribution of the player velocity in Fig. \ref{fig:histogram_nba} shows that our approach without theory-based computation had a wider distribution like ground truth than other models, and those with theory-based computation had less zero-velocity bins like ground truth than other models. }
\begin{table}[ht!]
\centering
\scalebox{1}{
\begin{tabular}{l|c}%|
\Xhline{3\arrayrulewidth} %\hline
% & \me{4}{carla simulation dataset}  \\ 
 & \me{1}{$L_{Covariates}$ at endpoint} \\
\hline
RNN & 1.637 $\pm$ 0.0096
\\DSW \cite{liu2020estimating}& ---
\\GCRN \cite{ma2021deconfounding} & ---
\\GCRN$+$X & 2.266 $\pm$ 0.0138\\
\hline
GV-CRN & 1.937 $\pm$ 0.0116
\\TG-CRN & \textbf{1.482} $\pm$ \textbf{0.0090}
\\TV-CRN & 1.566 $\pm$ 0.0094
\\ TGV-CRN (full) & \textbf{ 1.510} $\pm$ \textbf{0.0093}
\\
\Xhline{3\arrayrulewidth} % \hline
\end{tabular}
}
\vspace{3pt}
\caption{\label{tab:nba_end} \textcolor{black}{Endpoint prediction error on NBA dataset.}}
\end{table}

\begin{figure}[h]
\centering
\includegraphics[width=0.5\textwidth]{fig/histogram_nba.png}
\caption[]{\textcolor{black}{Histograms of player velocity in NBA dataset.}
}
%
\label{fig:histogram_nba}
\end{figure}

% \appendix
 

% Can use something like this to put references on a page
% by themselves when using endfloat and the captionsoff option.
\ifCLASSOPTIONcaptionsoff
  \newpage
\fi



% trigger a \newpage just before the given reference
% number - used to balance the columns on the last page
% adjust value as needed - may need to be readjusted if
% the document is modified later
%\IEEEtriggeratref{8}
% The "triggered" command can be changed if desired:
%\IEEEtriggercmd{\enlargethispage{-5in}}

% references section

% can use a bibliography generated by BibTeX as a .bbl file
% BibTeX documentation can be easily obtained at:
% http://mirror.ctan.org/biblio/bibtex/contrib/doc/
% The IEEEtran BibTeX style support page is at:
% http://www.michaelshell.org/tex/ieeetran/bibtex/
%\bibliographystyle{IEEEtran}
% argument is your BibTeX string definitions and bibliography database(s)
%\bibliography{IEEEabrv,../bib/paper}
%
% <OR> manually copy in the resultant .bbl file
% set second argument of \begin to the number of references
% (used to reserve space for the reference number labels box)

% biography section
% 
% If you have an EPS/PDF photo (graphicx package needed) extra braces are
% needed around the contents of the optional argument to biography to prevent
% the LaTeX parser from getting confused when it sees the complicated
% \includegraphics command within an optional argument. (You could create
% your own custom macro containing the \includegraphics command to make things
% simpler here.)
%\begin{IEEEbiography}[{\includegraphics[width=1in,height=1.25in,clip,keepaspectratio]{mshell}}]{Michael Shell}
% or if you just want to reserve a space for a photo:

% You can push biographies down or up by placing
% a \vfill before or after them. The appropriate
% use of \vfill depends on what kind of text is
% on the last page and whether or not the columns
% are being equalized.

%\vfill

% Can be used to pull up biographies so that the bottom of the last one
% is flush with the other column.
%\enlargethispage{-5in}



% that's all folks
\end{document}


